%% main.tex : Scientific Reports article, English manuscript (assembled 2026-10-05).
%% Template: Springer Nature LaTeX template v3.1 (December 2024), paper/manuscript/template/
%%   (zip sha256 812e76dcaa9c28dc1bff1fb6065d51729b67d4ea140552a05088317414a3ecae).
%% Build: ./build.sh (xelatex, bibtex with sn-nature.bst, three passes; logs in build/).
%% Engine: the 'pdflatex' class option is kept under xelatex. Without it the class loads breakurl,
%%   which fails under xelatex (docs/research/2026-10-04/scirep_format_and_drafts.md 5.2, test 4).
%%   Placeholders still contain Korean text, so xeCJK is loaded under XeTeX; a pdflatex build
%%   works only after all [DECISION/미확인/MISSING/X?] placeholders are replaced.
%% Order (MANUSCRIPT_SPEC 1.1, Sci Rep): title page and abstract, Introduction, Results,
%%   Discussion, Methods (ends with Code availability), Data availability, References,
%%   Acknowledgements, Author contributions, Additional information, Figure legends,
%%   Table 1, Figures 1-7 (PDF only).
%% Submission: the template asks for one .tex file without \input; flatten with latexpand
%%   (MANUSCRIPT_SPEC / scirep_format_and_drafts.md 5.4) before upload.
\documentclass[pdflatex,sn-nature]{sn-jnl}% Style for submissions to Nature Portfolio journals

\usepackage{graphicx}%
\usepackage{multirow}%
\usepackage{amsmath,amssymb,amsfonts}%
\usepackage{amsthm}%
\usepackage{mathrsfs}%
\usepackage[title]{appendix}%
\usepackage{xcolor}%
\usepackage{textcomp}%
\usepackage{manyfoot}%
\usepackage{booktabs}%
\usepackage{catchfile}%
\usepackage{iftex}%
\ifXeTeX
  \usepackage{xeCJK}% Korean text inside placeholders only
  \setCJKmainfont{Noto Serif CJK KR}%
  \setCJKsansfont{Noto Sans CJK KR}%
  \setCJKmonofont{Noto Sans CJK KR}%
  \xeCJKsetup{CJKspace=true}% keeps the spaces between Korean words (scirep_format_and_drafts.md 5.2, test 5)
  \xeCJKDeclareCharClass{CJK}{"0370 -> "03FF}% Greek letters typed inside placeholders (math Greek is unaffected)
\fi

\raggedbottom
\unnumbered% Scientific Reports headings are not numbered
\setcitestyle{super,open={},close={}}% Nature-style superscript citations; the sections place \cite directly after the word

% Figures and Table 1 body (read only; generated by scripts/4_visualization/paper_v3/)
\newcommand{\figdir}{../../../outputs/figures/paper/v3_restructure}

% Table 1: take the tabular from Table1_data.tex; the title and notes are in sections/legends.tex
% (Table caption), so the generated caption and notes paragraph of Table1_data.tex are not repeated.
\CatchFileDef{\TableOneRaw}{\figdir/Table1_data.tex}{}
\long\def\ExtractTableOne#1\noindent\begin#2\bottomrule#3\EndExtractTableOne{%
  \def\TableOneTabular{\begin#2\bottomrule\end{tabular}}}
\expandafter\ExtractTableOne\TableOneRaw\EndExtractTableOne

% One figure per page at the drawn width (170 mm), centred on the text block.
% A missing PDF is replaced by a framed box so that the build does not stop.
\newcommand{\PlaceFigure}[2]{%
  \begin{figure}[p]
  \centering
  \IfFileExists{\figdir/#2}{%
    \makebox[\textwidth][c]{\includegraphics[width=170mm,height=0.9\textheight,keepaspectratio]{\figdir/#2}}%
  }{%
    \fbox{\parbox[c][80mm][c]{0.9\textwidth}{\centering\large Figure #1 in preparation}}%
  }%
  \par\medskip
  {\small\bfseries Figure #1}
  \end{figure}
  \clearpage}

\begin{document}

% Block 1 of front.tex: title, authors, abstract, keywords
\def\frontblock{1}% front.tex : 영문 원고의 제목 쪽, 초록, 핵심어, Code availability, Data availability,
%             Acknowledgements, Author contributions, Additional information(Competing interests)
% 작성 2026-10-04. LaTeX 본문 조각이다(프리앰블 없음). 명세: paper/manuscript/MANUSCRIPT_SPEC.md
%   1.1(순서), 1.3(제목·핵심어), 1.5(서술 방향), 2.1-2.4(초록), 6.15(M15 Code availability),
%   6.16(Data availability), 13.1(사용자 결정 대기).
% 서식: Springer Nature 템플릿 sn-jnl(sn-nature 선택지)의 명령(\title, \author, \affil, \abstract,
%   \keywords, \bmhead)을 쓴다. 블록 1-4 는 원고 조립 때 아래 표시한 자리로 옮긴다.
%   블록 1: \begin{document} 뒤, \maketitle 앞
%   블록 2: Methods 마지막 소절(M15). methods.tex 에 같은 소절이 있으면 하나만 남긴다.
%           M14 Use of large language models 는 methods.tex 의 몫이고 문구는
%           [DECISION: LLM 사용 기재 문구] 다(이 파일에는 본문을 쓰지 않았다).
%   블록 3: 본문 끝, \bibliography 앞(Sci Rep 점검표: Data availability 는 References 앞)
%   블록 4: \bibliography 뒤(\backmatter). Acknowledgements, Author contributions, Additional information
% 인용 키: paper/manuscript/en/refs_front.bib. 키 이름은 다른 절의 FirstAuthorYear 규칙을 따르고,
%   Streletskiy2025 와 munozsabater2021_era5_land 는 refs_intro.bib, refs_results_b.bib 와 같은 키다.
%
% 수치 출처(원고 조립 때 shared/numbers.tex 매크로로 바꾼다. 자릿수는 지침 4.5):
%   -1.6 cm, 95% 신뢰구간 -2.0 to -1.0 .... C1 README 2.3 AB3(-1.63 [-1.98, -1.04], Holm P 0.001)
%   -2.5 cm, -3.0 to -1.9, 1지역 ......... C2 README 2.1 A6(AB4 -2.45 [-3.04, -1.88]), 지역 행 C3 README E36a
%   잔차 ML 추가 차이 확인 못함 ............ C2 README 2.1 A7(AB5 -0.18, Holm P 0.136, 초록 규칙 (b))
%   -0.87, -1.3 cm, 2지역, 캐나다 ......... C5 README 2.1 E2, E3, 2.2 E6, E7(WF4-a, 결과 열람 뒤 설계)
%   13.9, 14.4 cm(라벨 1000개) ............. C4 README 2.1 E1.R1.n1000(13.87 / 14.41)
%   약 99% ................................ C8 README 2.2 B4(격자 안 몫 0.9%, 알래스카 지역 내, 라벨 전량)
%   -2.5, +1.7 cm(라벨 100개) .............. C6 README 2.1 A3, A17(WF2-a 서술 행)
%   다섯 지역, 100 km, 라벨 0-전량 ......... D README 1.3, E24, E25
%   원천 CSV 로 다시 확인한 값(2026-10-04): lgw_bundle.csv(AB3 -1.629003, AB4 -2.454509, AB5 -0.177787),
%   wf_tests.csv(WF4-a -0.873044, -1.320595; WF2-a -2.538178, +1.651555), wf2b_tests.csv(13.86715, 14.40761),
%   wf3b_decomp.csv(share_gain_within 0.008612).
%
% 이 파일의 자리표시: [DECISION: ...] 사용자 결정, [미확인: ...] 확인하지 못한 사실,
%   [XC: ...] [XE: ...] 새 실험(명세 11절 등록부의 문자열 그대로). 제출 전 점검에서 0건이어야 한다.
% 빌드 주의: 자리표시에 한글이 있어 자리표시가 남은 동안은 xelatex(kotex)로 조립한다.
%   pdflatex 빌드는 자리표시를 모두 바꾼 뒤에 한다(지침 M-11).

% 조립(2026-10-05): main.tex 가 \frontblock 을 1, 3, 4 로 정한 뒤 이 파일을 블록마다 \input 한다.
%   각 블록은 \ifnum\frontblock=N ... \fi 로 감쌌다(0 이면 아무것도 내지 않는다). 블록 2 는
%   methods.tex 의 Code availability(M15)와 같은 소절이라 main.tex 가 부르지 않는다(위 블록 2 지시: 하나만 남긴다).
\providecommand{\frontblock}{0}

% =====================================================================
% 블록 1. 제목 쪽, 초록, 핵심어
% =====================================================================
\ifnum\frontblock=1\relax

% 임시 제목(사용자 결정 2026-10-04: 대회 제목의 영역, 나중에 바꾼다). 14단어, 콜론·의문문·금지어 없음.
% [DECISION: 최종 제목]
\title{Permafrost active-layer thickness prediction under sparse observations using physics-based pseudo-label augmentation and machine learning}

% 저자 후보 기록(명세 13.1 의 3번): 제1 저자 후보 백승원(Seungwon Baek), KOPRI 공동 연구 가능자
% 김승희(Seung Hee Kim), 정윤택(Yoon Taek Jung). 저자 자격은 이 파일이 정하지 않는다.
% 형식(템플릿): \author*[1]{\fnm{이름} \sur{성}}\email{...}, \affil*[1]{\orgdiv{}, \orgname{}, \orgaddress{...}}
\author*[1]{\fnm{[DECISION: 저자, 소속, 교신 저자]}}\email{[DECISION: 교신 저자 전자우편]}

\affil*[1]{\orgname{[DECISION: 저자, 소속, 교신 저자]}}

% 초록: 판 B(사용자 잠정 선택, WF 결과는 판정어 없는 수치). 명세 2.2 문안에서 바꾼 곳은 셋이다.
%   (1) 약어 첫 등장 풀이(지침 M-07): RMSE 를 S3 에서 풀어 쓰고, '95% CI' 를 '95% confidence interval' 로 썼다.
%   (2) WF 결과 앞에 '결과 열람 뒤 설계' 표지를 넣었다(명세 1.5, 지침 4.5 끝 항목).
%   (3) 200단어를 지키려고 S2('ML models ... without recalibrated baselines'), S3('Stefan models'),
%       S7(지역을 괄호로), S8('Sub-grid covariates are needed next.')을 줄였다.
% 단어 수(wc -w, 자리표시 제외): 199. 판정어는 S4(AB3 규칙 (a), AB4 규칙 (a)와 지역 수, AB5 규칙 (b))에만 있다.
% [DECISION: 초록 판 A 또는 B]  [DECISION: 판 B 의 WF 수치 개수 허용 여부]
% S8 의 두 자리표시는 결과가 나오면 FINAL_BATCH 의 해당 문장으로 바꾸고, XC 수치를 넣으면 S7 이나 S8 을 대신해 200단어를 지킨다.
\abstract{We provide a label-count workflow for mapping active-layer thickness in sparsely measured permafrost regions with the Stefan model and machine learning (ML). ML models are mostly tested within training regions at one label count, without recalibrated baselines. We held out five regions with 100~km buffers and compared root-mean-square error (RMSE) from zero to all labels against Stefan models with source and recalibrated coefficients (internally pre-registered). Without labels, Stefan pseudo-labels gave lower error than shuffled ones ($-1.6$~cm; 95\% confidence interval $-2.0$ to $-1.0$); with ten labels, recalibration gave lower four-region mean error than the source coefficient ($-2.5$~cm; $-3.0$ to $-1.9$; lower in one region), with no further difference established for residual ML. In analyses designed after these results, method selection by within-target cross-validation with 40 and 160 labels changed RMSE by $-0.87$ and $-1.3$~cm against a fixed residual recipe (two regions, change from Canada). Within Alaska, residual ML reached 13.9 against 14.4~cm RMSE for recalibrated Stefan (1000 labels); about 99\% of the squared-error reduction (all labels) lay between climate grid cells. Spreading 100 labels over blocks changed RMSE against random sampling by $-2.5$~cm (Canada) and $+1.7$~cm (Lena Delta). Sub-grid covariates are needed next. \textcolor{blue}{[PENDING: XE-a, XE-b, satellite inputs (stage 2); registered deadline 8 October 2026, 14:22 KST]}}
% [DECISION: abstract XC sentence] XC is left out of the abstract (FINAL_BATCH 7.1 allows numbers only or exclusion; a numbers-only
% sentence replaces S7 (19 words) and makes the abstract 214 words). Candidate (34 words; would need 14 further words cut):
%   In new random splits of the same regions, the workflow changed RMSE by -1.51 cm against recalibrated Stefan at 160 labels
%   (pooled across regions; over 0.5 cm larger error in 2 of 3 regions).

% 비교판(판 A, WF 제외). 명세 2.3 문안에 같은 약어 수정('95% confidence interval')만 했다. 197단어.
% 판 A 를 고르면 위 \abstract{} 의 내용을 아래 문단으로 바꾼다. 판정어 근거는 명세 2.3 끝 문단.
%   We provide a label-count workflow for mapping active-layer thickness (ALT) in sparsely measured
%   permafrost regions with the Stefan thaw model and machine learning (ML). ML models of ALT are usually
%   evaluated within training regions, at one label count and without a recalibrated physical baseline.
%   We held out four Arctic regions with 100~km buffers and compared errors from zero to all target labels
%   against the Stefan model with source and recalibrated coefficients, in an internally pre-registered
%   analysis. Without labels, Stefan pseudo-labels gave lower error than shuffled ones ($-1.6$~cm; 95\%
%   confidence interval $-2.0$ to $-1.0$), and no difference was established between a physics-model
%   ensemble and the source-coefficient model. With ten labels, coefficient recalibration gave lower
%   four-region mean error ($-2.5$~cm; $-3.0$ to $-1.9$; one of four regions), no further difference was
%   established for residual ML, and residual learning gave lower error than physics outputs used as ML
%   inputs. With all labels, residual ML gave lower four-region mean error than the source-coefficient
%   model, mainly from one region, and lower error than the recalibrated model by less than 0.5~cm. The
%   workflow therefore uses physics at every label count and assigns ML a small correction once labels
%   allow recalibration.

% 핵심어 6개(Sci Rep 상한 6). 명세 1.3 제안. [DECISION: 핵심어]
\keywords{Active-layer thickness, Permafrost, Stefan model, Physics-guided machine learning, Spatial transfer, Sampling design}

\fi

% =====================================================================
% 블록 2. Code availability(Methods 마지막 소절, 명세 6.15 M15)
% =====================================================================
\ifnum\frontblock=2\relax
% 코드 서체(\texttt)는 이 소절의 실제 저장소 경로에만 쓴다(지침 4.1, M-10).
% 경로는 2026-10-04 에 저장소에서 확인했다. 패키지 판은 FINAL_BATCH 1절 '패키지 판' 행,
%   이전 실행의 판은 방법 초안 Table 4(Supplementary Table S16 으로 옮길 예정).
% 저장소 후보: git 원격 git@github.com:01willy/Polar_Bigdata(현재 공개 여부와 공개 시점은 결정 대기).
% TabPFN 판 8.0.7 의 Prior Labs License(v1.2) 10항은 가중치나 원본을 배포할 때 'Built with PriorLabs-TabPFN'
%   표시를 요구한다. 가중치를 재배포하지 않으므로 본문에는 넣지 않았다. 저장소 README 에 넣을지는 공개 때 판단한다.
\subsection*{Code availability}

The code for data assembly, experiments, aggregation, figures and tests is available at [DECISION: 공개 저장소 주소와 공개 시점] and will be archived under the release tag of the reported analyses at [DECISION: Zenodo DOI]. The experiment scripts under \texttt{scripts/} are \texttt{h40\_label\_grid.py} and \texttt{h42\_label\_grid\_ext.py} (transfer), \texttt{h41\_validation\_ladder.py} (validation ladder), \texttt{h54\_workflow.py} (within-region tests), \texttt{h39\_scenarios.py} (aggregation), \texttt{wf0\_reanalysis.py} (post hoc reanalysis) and \texttt{x\_*.py} (additional analyses). The repository also holds the unit tests (\texttt{tests/}), the cloud job configurations (\texttt{configs/rescale/}) and the registration documents with their revision histories. Each output file of a run records the code hash, the configuration hash and the package versions, and the aggregate metadata record the git commit. The cloud runs of the additional analyses use CatBoost 1.2.10, scikit-learn 1.9.1, pandas 2.3.3, SciPy 1.17.1 and NumPy 2.1.1; the versions of earlier runs are listed in Supplementary Table S16. The pretrained weights of TabPFN v2 and TabICL v2 are not redistributed and are obtained from their providers.

\fi

% =====================================================================
% 블록 3. Data availability(본문 끝, References 앞, 명세 6.16)
% =====================================================================
\ifnum\frontblock=3\relax
% 확인 기록(2026-10-04):
%   - 자료 인용 19건은 방법 초안 References 표(67행)와 references/INDEX.md 12.4 표에 있고,
%     DOI 와 제목·연도·저장소를 DataCite API 로 다시 대조했다(refs_front.bib 각 항목 주석).
%   - 새 지역 자료원의 약관은 data/processed/fidelity_base_v4_meta.json 의 by_source 를 따랐다.
%     약관 확인분(license_attention False): calm_pangaea_v4add, firealt, Disko, Ilulissat, myThaw,
%     Palmtag pedon, Du(Zenodo, MIT), Fu(figshare, 지온 유도, 라벨 정의 시험), Kolyma, Indigirka,
%     Mamontov Klyk, Syrdakh, Tavvavuoma, Yamal(PANGAEA 행). 약관 미확인: CALM 누리집 하위 지점(Abisko,
%     Kapp Linné), CUSP v1.1, NSIDC GGD353, Yamal 보고서 유래 행, NSIDC GGD402(기술 통계만).
%     약관 미확인 자료의 인용(Åkerman 1998, Åkerman and Johansson 2008, Jorgenson and Kanevskiy 2025,
%     Petrone 2016, Nixon 2003)은 Supplementary Table S17 과 SI 참고문헌에 둔다. [DECISION: 참고문헌 처리]
%   - CCI ALT v4.0 의 DOI 는 DataCite 에서 제목 'Permafrost active layer thickness for the Northern
%     Hemisphere, v4.0'(2024, CEDA)으로 확인했다. INDEX 01절 esa_cci_permafrost_v4 의 DOI(7479606004...)는
%     DataCite 에 없다. 명세 9.3 의 불일치는 ALT 자료 인용 쪽에서는 d34330ce... 로 정리된다.
%   - SAR: 10.3334/ORNLDAAC/2004 는 DataCite 제목이 'ABoVE: Active Layer Thickness from Airborne L- and
%     P- band SAR, Alaska, 2017, Ver. 3'(Chen et al., 2022)이고, 내려받기 스크립트
%     scripts/0_download/resalt_insar.py 가 이 DOI 로 받은 파일(PDO_ReSALT_*_2017_03.nc4)이 명세의 'ReSALT' 다.
%     PolSAR 상향 ALT(data/raw/polsar_alt/upscaled_alt_{2014,2015,2017}.tif)의 출처는 내려받기 기록이 없다.
%     DataCite 10.3334/ORNLDAAC/2332 의 제목('ABoVE: Upscaled Active Layer Thickness in Northern Alaska,
%     2014-2017', Whitcomb et al., 2025)이 같은 자료로 보이나 확인하지 못했다.
%   - Copernicus DEM DOI 10.5270/ESA-c5d3d65 는 doi.org 에서 COP-DEM 자료 설명 쪽으로 연결된다(인용 형식은 미확인).
%   - 지도 마스크(CCI 영구동토 범위 평균, JRC 지표수)는 방법 초안 Data availability 의 문장을 따랐다.
\section*{Data availability}

The active-layer labels of the main regions are available from the GTN-P/CALM compilation\cite{Streletskiy2025} (PANGAEA, \url{https://doi.org/10.1594/PANGAEA.972777}), the ALLena compilation for the Lena River Delta\cite{Veremeeva2025} (PANGAEA, \url{https://doi.org/10.1594/PANGAEA.973813}) and the ABoVE soil moisture and active-layer thickness data set, version~2\cite{Moore2025} (ORNL DAAC, \url{https://doi.org/10.3334/ORNLDAAC/2369}). The labels of the additional regions come from published data sets\cite{DuE2026,Grosse2007,Palmtag2022,Talucci2024,Pohl2026,Makarieva2017,Liang2023,Walker2009,Scheer2024,Martin2023,Boike2024,Hammar2025,Zastruzny2024,Sannel2020}, and the temperature-derived labels used only in the label-definition test come from a Tibetan Plateau compilation\cite{Fu2025}. Their identifiers and licences are listed in Supplementary Table S17. Labels from sources without a confirmed redistribution licence were excluded from the licence-verified edition used for the main results and are not redistributed; they can be obtained from the original providers listed in Supplementary Table S17. The full North Atlantic edition, which includes such labels, is reported in the Supplementary Information only.

The covariates were derived from ERA5-Land monthly averaged data\cite{munozsabater2021_era5_land} (Copernicus Climate Data Store, \url{https://doi.org/10.24381/cds.68d2bb30}), SoilGrids 2.0\cite{Poggio2021} [미확인: 접근 시점의 SoilGrids 판 표지], the Copernicus DEM GLO-30 (\url{https://doi.org/10.5270/ESA-c5d3d65}) [미확인: Copernicus DEM 인용 형식] and the ESA CCI Permafrost active-layer thickness product, version 4.0\cite{Westermann2024} (CEDA, \url{https://doi.org/10.5285/d34330ce3f604e368c06d76de1987ce5}). The synthetic-aperture radar (SAR) covariates of the within-Alaska test came from the ABoVE airborne L- and P-band SAR active-layer retrievals for 2017 (ORNL DAAC, \url{https://doi.org/10.3334/ORNLDAAC/2004}) and from upscaled P-band polarimetric SAR active-layer maps of northern Alaska\cite{Whitcomb2024} [미확인: PolSAR 상향 ALT 자료의 DOI, 후보 10.3334/ORNLDAAC/2332]. The map masks use the ESA CCI Permafrost extent product [미확인: 영구동토 범위 자료의 DOI와 판] and the JRC Global Surface Water occurrence layer\cite{Pekel2016}.

The assembled label and covariate tables, the run outputs, the aggregate and verdict tables, the cell-level predictions and the source data of all figures will be deposited at [DECISION: 저장소와 DOI]. Rows from sources without a confirmed licence are excluded from the deposited tables. [DECISION: KPDC 자료, SI 포함 여부와 식별자, 이용 조건]

\fi

% =====================================================================
% 블록 4. \bibliography 뒤(\backmatter)
% =====================================================================
\ifnum\frontblock=4\relax
% Sci Rep: 감사의 글은 짧게, 저자 기여는 참고문헌 뒤에 모든 저자별로, 이해상충은 'Additional information' 아래.
% 세 문구 모두 [DECISION: Acknowledgements, Author contributions, Competing interests 문구](명세 13.1 의 5번).

\bmhead{Acknowledgements}

[DECISION: Acknowledgements, Author contributions, Competing interests 문구] We thank the providers of the data sets listed under Data availability. The active-layer thickness covariate uses data of the ESA Climate Change Initiative Permafrost\_cci project. Terrain covariates were produced using Copernicus WorldDEM-30 \textcopyright{} DLR e.V. 2010--2014 and \textcopyright{} Airbus Defence and Space GmbH 2014--2018, provided under COPERNICUS by the European Union and ESA; all rights reserved. [미확인: Copernicus DEM 사용 허가가 요구하는 귀속 문구의 원문 대조] [미확인: ERA5-Land(Copernicus Climate Change Service) 사용 허가가 요구하는 귀속 문구] [DECISION: KPDC 자료, SI 포함 여부와 식별자, 이용 조건] [DECISION: 대회 보고서 인용 여부] [DECISION: 연구비 지원 기관과 과제 번호, 또는 지원 없음]

% KPDC 자료를 SI 에 두면 KPDC 감사 문구를 더한다(앞부분 초안 3절의 문안. 식별자 2707, 2955 확인 뒤).
% 대회: '2026 극지 빅데이터-인공지능 활용 경진대회'(주최 극지연구소, 주관 KPDC). 공식 영문 명칭 [미확인].
%   선행 공개 사실은 표지 편지에 적고 사본을 첨부한다(명세 13.1 의 12번).

\bmhead{Author contributions}

[DECISION: Acknowledgements, Author contributions, Competing interests 문구] [DECISION: 저자, 소속, 교신 저자]

% 형식(결정 뒤): 모든 저자에 대해 CRediT 역할(Conceptualization, Methodology, Software, Validation,
%   Formal analysis, Investigation, Resources, Data curation, Writing (original draft), Writing (review and
%   editing), Visualization, Supervision, Project administration, Funding acquisition)을 이름과 함께 적고,
%   끝에 'All authors reviewed the manuscript.' 같은 문장을 둔다. 저자 자격은 이 파일이 정하지 않는다.

\bmhead{Additional information}

Competing interests: [DECISION: Acknowledgements, Author contributions, Competing interests 문구]

Correspondence and requests for materials should be addressed to [DECISION: 저자, 소속, 교신 저자].

% Competing interests 문안(결정 뒤, 해당 없을 때 Sci Rep 예시 문구): 'The authors declare no competing interests.'
\fi


\maketitle

\section{Introduction}\label{sec:intro}
% intro.tex : Introduction (English manuscript, body text only, no preamble)
% 작성 2026-10-04. 명세: paper/manuscript/MANUSCRIPT_SPEC.md 3절(I1-I4, 700단어 목표).
% 인용 키: paper/manuscript/en/refs_intro.bib (refs_discussion.bib 와 같은 FirstAuthorYear 규칙).
%
% 수치 출처(원고 조립 때 shared/numbers.tex 매크로로 바꾼다. 자릿수는 지침 4.5):
%   30 / 13,606 / 343 ............ paper/claims/D_data_and_design/README.md 2절 E14, E10
%   1.31 / 2.51 / 11.03 .......... 같은 README E38 (fig1_z_summary.csv E_mean, 셀 E 기하 평균)
%   1.6 [1.0, 2.0], P = 0.001 .... C1_label0_safety/README.md 2.3 첫 행(셀 가중 -1.63 [-1.98, -1.04])
%   2.5 [1.9, 3.0], P = 0.001 .... C2_bias_diagnosis/README.md 2.1 A6(-2.45 [-3.04, -1.88]), 지역 행 1/4 은 C3 README E36a
%   거의 전부(약 99%) ............. C8_gain_source/README.md 1.3, B4(알래스카 지역 내, 결과 열람 뒤 설계)
%   0.29 degC(2007-2016) ......... Biskaborn 2019 초록(로컬 PDF 대조)
%   68 지점, 무작위 70/30 1회, 시험 R2 0.24 대 0.54 ... Gautam 2025 본문(로컬 PDF 대조)
%   35 years ..................... Streletskiy 2025 자료 제목(DataCite 대조)
% 열린 결정: 지역 영문 표기(W Russia, E Russia, Lena Delta)는 MANUSCRIPT_SPEC 13.1 의 7번 결정 대기.
% 새 실험(XA-XJ) 자리표시: 명세 11절 등록부에 서론 자리가 없어 넣지 않았다.

Active-layer thickness (ALT) is the maximum seasonal thaw depth of the ground above permafrost, and its change affects infrastructure built on permafrost\cite{Karjalainen2019,Ran2022b}. Permafrost temperature increased by 0.29\,$^{\circ}$C on global average between 2007 and 2016\cite{Biskaborn2019}. ALT is measured at a limited number of sites, many of them in the 35-year Circumpolar Active Layer Monitoring (CALM) compilation\cite{Streletskiy2025}. In the five regions evaluated here, the number of measured ALT values (labels) ranges from 30 in E Russia to 13,606 in Alaska, and the Alaskan labels lie at 343 locations of 1\,km (Table~1). The Stefan solution\cite{Stefan1891} predicts ALT as a coefficient $E$ times the square root of the thawing degree-days (TDD), with $E$ calibrated by local thaw-depth measurements\cite{Nelson1997}. The geometric mean of $E$ over labelled cells is 1.31\,cm\,($^{\circ}$C\,d)$^{-1/2}$ in the Lena Delta, 2.51 in W Russia and 11.03 on the Tibetan Plateau (Fig.~1b). However, because measurements in a new permafrost region (the target region) are few, the error of ALT maps there depends on models fitted to other regions (the source data).

Outputs of the Stefan or Kudryavtsev models have served as machine-learning (ML) input features\cite{Pilyugina2025,WangG2025}, $E$ has been predicted from covariates\cite{Ran2022b,ZhangC2024}, labels generated by physical rules have trained a zero-curtain detector\cite{Gay2026}, and process-model pre-training and residual modelling are used in other environmental fields\cite{Read2019,Jia2021,Willard2022}. Each component has precedents, and the studies differ in how they were evaluated. Gautam et al. compared a random forest with a Stefan model at 68 CALM sites in Alaska using one random 70/30 split and reported test coefficients of determination of 0.24 and 0.54, respectively\cite{Gautam2025}. Large-scale ALT maps have been validated with distance blocks\cite{Aalto2018,Karjalainen2019,Ran2022a}. For clustered spatial data, random cross-validation can overstate map accuracy\cite{Roberts2017,Ploton2020,MeyerPebesma2022}, whereas Wadoux et al. argued that unbiased map accuracy requires probability sampling with design-based inference rather than spatial cross-validation\cite{Wadoux2021}. The ALT evaluations above each use one label condition, whereas adjacent fields have evaluated models in held-out regions or with few local observations, in some cases against calibrated process models, for lake temperature\cite{Read2019,Jia2021,Willard2021}, ungauged catchments\cite{Pool2019,Feng2023}, subsurface temperature\cite{OMalley2026} and spatial risk prediction\cite{Portes2026}. To our knowledge, no study of permafrost ALT or mean annual ground temperature has reported the error in regions excluded from the training data as a function of the number of target labels, relative to a Stefan model recalibrated with those labels.

A new region needs a rule for combining the Stefan model and ML as target labels accumulate, and a procedure for adding labels. We address three questions. (1) When no target labels are available, what does physics information add to ML, and where does the gain come from? (2) As target labels increase from ten to thousands, how much does ML add to a Stefan model recalibrated with the same labels, and at which spatial scale does the gain arise? (3) As labels accumulate, how should the method be selected, and by which procedure should new labels be placed? The answers define the steps of a workflow.

Each target region was held out from the source data with a 100\,km buffer, labels were drawn from half of its 0.5$^{\circ}$ blocks, and errors were scored on the other half, for 0, 3, 10, 40, 160, 320, 1000 and all available labels (Fig.~1c,d). Every method was compared with the Stefan model before and after recalibration on the same labels, as in lake-temperature studies\cite{Read2019,Jia2021}, with the recalibrated coefficient shrunk towards the source value\cite{Steyerberg2004}. Within label-rich regions, block-holdout tests designed after earlier results were viewed measured the gain of ML over the recalibrated model. Contrasts were judged from block-bootstrap 95\% confidence intervals (CIs) with a $\pm$0.5\,cm equivalence margin, and the transfer analysis was internally pre-registered. We find that, without target labels, Stefan pseudo-labels (Stefan predictions used as training labels) reduced the mean root-mean-square error (RMSE) of four main regions (Lena Delta, Canada, W Russia and E Russia) relative to shuffled pseudo-labels by 1.6\,cm (95\% CI 1.0 to 2.0; Holm-adjusted $P = 0.001$), and that recalibrating the Stefan coefficient with ten labels reduced it by 2.5\,cm relative to the source coefficient (95\% CI 1.9 to 3.0; Holm-adjusted $P = 0.001$; lower error in one of the four regions). Within Alaska, with all labels, almost all of the gain of ML over the recalibrated model came from correcting bias between climate grid cells. The contribution is the evaluation design, its results including negative ones, and the label-count workflow built from them.


% results_a.tex : Results, first four subsections (evaluation design, zero labels,
% three to ten labels, method selection). Body text only, no preamble.
% Spec: paper/manuscript/MANUSCRIPT_SPEC.md sections 4.1 to 4.4 and 11.
% Numbers: copied from paper/claims/{D_data_and_design,C7,C1,C2,C3,C5}/README.md and
% rounded with style guide 4.5 (|delta| < 1 cm two decimals, other cm values one decimal,
% correlations two decimals, percentages integer, CI ends with the digits of the printed
% point estimate). Rounding used the full-precision source rows that the READMEs cite.
% Replace with numbers.tex macros when build_numbers.py exists (QA item M-09).
% Sign: negative values use U+2212. Intervals are cell-weighted unless marked block-equal.
% Bracketed new-experiment placeholders are the registered strings of spec section 11 and
% must be replaced only by the pre-fixed interpretation sentences. The two bracketed
% missing-value markers follow spec section 12, item M-04 (median and mean without W Russia).
% Placeholders contain Korean text, so a build that keeps them needs a Korean font setup.
% Citation keys follow refs_intro.bib and refs_discussion.bib (FirstAuthorYear) and are
% listed in paper/manuscript/en/refs_results_a.bib.
% This file opens the Results section; results_b.tex continues with subsections only.

\section{Results}\label{sec:results}

\subsection{Evaluation design and validation ladder}\label{sec:results-design}

All methods were compared on one protocol in which each target region and a 100~km buffer were removed from the source data, labels were drawn from half of the target's 0.5° blocks, and errors were scored on the other half, for label counts from zero to all available rows (30 to 13,606 rows per region) (Fig.~1c,d; Table~1).
Every method was compared with two Stefan baselines that used the same target labels: the source-coefficient Stefan model, whose coefficient $E_0$ was estimated from the source cells (1.59--1.62~cm~(°C~d)$^{-1/2}$ for the main regions, whose source cells were 78--94\% Alaskan, and 1.50 for the Alaskan reference target), and the Stefan model recalibrated with the target labels by shrinkage towards $E_0$ ($\kappa = 10$).
Four-region means weight the Lena Delta, Canada, W Russia and E Russia equally and exclude Alaska, which is a held-out reference target.
The transfer comparisons were internally pre-registered, intervals are cell-weighted unless marked block-equal, and error floors used to scale effect sizes were 11.3, 14.8 and 17.9~cm in Alaska, the Lena Delta and Canada (Methods).

Averaged over five regions with equal weight, the RMSE of direct ML increased from 22.9~cm under random cell splits to 33.3~cm under region holdout (difference 10.3~cm; 95\% confidence interval (CI) 7.3 to 12.8; block-equal 6.1~cm, 4.7 to 7.6), whereas a Stefan model with the coefficient fitted by least squares to the training cells stayed between 26.8 and 27.0~cm (Fig.~1e).
The 10.3~cm increase is reported as an effect size rather than a confirmatory result, and the registered contrast between cluster holdout with a 500~km buffer and random cells was 8.2~cm (95\% CI 5.6 to 10.4; Supplementary Table S2).
Within Alaska, the median distance from test cells to the nearest training cell was 0.004~km under random splits and 53.5~km under k-fold nearest-neighbour distance matching (kNNDM) cross-validation, which mimics the 81.6~km from map cells to the nearest label\cite{Linnenbrink2024}; the RMSE of direct ML rose from 11.6 to 17.9~cm, and that of the Stefan model stayed at 14.2 to 14.5~cm (means of three seeds).
Random validation of clustered samples can overstate model accuracy\cite{Ploton2020}, and model-free, unbiased estimates of map accuracy have been argued to require probability sampling\cite{Wadoux2021}.
Our labels are a clustered non-probability sample, so the errors below are prediction errors on held-out blocks.
[XH: R1 | 결과 열람 뒤 설계, 서술(판정어 없음). L19, L20 은 등록 | 2.8 사전 고정 서술 규칙]

\subsection{Transfer without target labels}\label{sec:results-label0}

Without target labels, pseudo-labels from the Stefan model reduced RMSE relative to shuffled pseudo-labels by 1.6~cm (95\% CI 1.0 to 2.0; block-equal 0.6 to 1.5; four regions; Holm-adjusted \textit{P} = 0.001), and in a post hoc count, fewer of 30 targets lost more than 2~cm against the source-coefficient Stefan model with physics pseudo-label augmentation (2 targets) than with direct ML (18 targets) (Fig.~3b; Fig.~6a,b).
Relative to two constant placebos the reductions were 1.7 and 3.6~cm, and relative to a placebo linear in the square root of thawing degree-days 0.48~cm without labels (below 0.5~cm), with equivalence within ±0.5~cm at ten labels (registered verdict: mixed).
The 2~cm threshold was set after the transfer results had been viewed, and the 30 targets comprise seven regions and 23 sub-regions or label-expanded versions of them; counted by intervals above zero under both weightings, 14 targets had higher error with direct ML and 8 with physics pseudo-labels (post hoc computation; Supplementary Table S5).

For the five learners whose primary and common-resampling intervals agreed, direct ML had higher four-region mean error than the source-coefficient Stefan model: a random forest (2.6~cm) and four neural networks configured by leave-one-source-region-out cross-validation (2.4 to 5.2~cm) (Fig.~6c).
For the main learner, CatBoost, the difference was 2.2~cm (95\% CI 1.1 to 3.4; block-equal 1.6~cm, 0.5 to 2.7; Holm-adjusted \textit{P} = 0.023), and for two further CatBoost settings, TabPFN v2 and TabICL v2 it was 0.85 to 1.9~cm.
These five verdicts depend on the split-independence assumption, and the TabPFN v2 verdict was significant before correction only (Holm-adjusted \textit{P} = 0.117).
Direct CatBoost also had 3.2~cm higher error than the Stefan model driven by year-matched thawing degree-days (95\% CI 1.9 to 4.3; block-equal 2.0~cm, 0.9 to 3.1; auxiliary hypothesis with an unblinded anchor variant).
Residual models with a low weight ($\lambda = 0.25$) on the source-coefficient anchor were equivalent to that model within ±0.5~cm for CatBoost, TabPFN v2 and TabICL v2, a verdict that depends on the split-independence assumption.

Three results depart from this pattern.
With Alaska held out, physics pseudo-labels had 1.8~cm higher error than the source-coefficient model (95\% CI 0.4 to 2.6).
On the Tibetan Plateau, where the source coefficient is far from the regional value (Fig.~1b), direct ML had 60.0~cm lower error (95\% CI 57.1 to 62.9; an expected result from an unblinded, cross-environment comparison).
With ten labels, direct TabICL v2 had 2.9~cm lower four-region mean error (95\% CI 1.1 to 3.6; depends on the split-independence assumption) but 6.0~cm higher error in the three-region mean that includes Alaska.
[XG: R2 | 결과 열람 뒤 설계, 등록 이탈(WRAPUP 10) | 2.7 사전 고정 해석 문장 XG-1, XG-2, XG-3, 공통 문장]
We next asked what the first ten labels changed.

\subsection{Coefficient recalibration with three to ten labels}\label{sec:results-recal}

With ten target labels, recalibrating the Stefan coefficient reduced the four-region mean RMSE relative to the source coefficient by 2.5~cm (95\% CI 1.9 to 3.0; block-equal 1.9 to 3.3; Holm-adjusted \textit{P} = 0.001), with lower error in one of four regions (W Russia, 12.0~cm), higher error in Canada (2.3~cm; 95\% CI 0.7 to 2.9) and no difference established in the Lena Delta and E Russia (Fig.~2; Fig.~3a).
[MISSING: 중앙값과 러시아 서부 제외 값, build\_numbers.py 에서 lgw\_bundle.csv 지역 행으로 생성]
Adding residual ML to the recalibrated anchor changed RMSE by a further −0.18~cm (95\% CI −0.53 to −0.01; Holm-adjusted \textit{P} = 0.136), a difference significant before correction only, below 0.5~cm and dependent on the split-independence assumption.
In W Russia, 13.5~cm of the 14.0~cm reduction relative to the source coefficient with all labels came from recalibration (post hoc decomposition).

The error reduction of label-based correction (recalibration plus residual) was rank-correlated with the source-coefficient bias measured with ten labels (Spearman $\rho = 0.38$; 95\% CI 0.17 to 0.55; 28 targets; one-sided \textit{P} = 0.0002; designed after the label-grid results had been viewed), where the reduction is that of the method selected within the target labels (next section) relative to the source-coefficient model with all labels.
The sign of the bias was not related to the reduction ($\rho = -0.01$; 95\% CI −0.17 to 0.18) (Fig.~4a,b; Supplementary Table S8).
We could not establish a relation between coefficient error and the gain of ML beyond recalibration (post hoc rank correlations −0.08 to 0.39 depending on the definition; Supplementary Table S13).
[XA: R3 | 사후 분석(재현(비맹검)) | 2.1 해석 조각]

With ten or fewer labels, the residual structure had lower four-region mean error than ML that receives Stefan outputs as inputs (physics-input ML), by 2.7~cm without labels (95\% CI 1.7 to 4.3) and 3.7~cm with ten labels (95\% CI 2.8 to 5.2; Holm-adjusted \textit{P} = 0.001) (Fig.~3c).
[MISSING: 중앙값과 러시아 서부 제외 값, build\_numbers.py 에서 lgw\_bundle.csv 지역 행으로 생성]
The physics-input model of the ten-label contrast does not receive the recalibrated coefficient; in the contrast that supplies recalibrated Stefan values as inputs, the residual structure still had 2.5~cm lower error (95\% CI 1.6 to 3.3), with lower error in W and E Russia, 3.1~cm higher error in Canada (95\% CI 1.4 to 4.0; the E Russia and Canada verdicts depend on the split-independence assumption) and no difference established in the Lena Delta.

Adding physics pseudo-labels to residual ML was equivalent within ±0.5~cm with ten labels (−0.13~cm; 95\% CI −0.31 to 0.22; Holm-adjusted equivalence \textit{P} = 0.002) and increased error with 40 labels in the Lena Delta and Canada by 0.34~cm (95\% CI 0.25 to 0.49), a difference below 0.5~cm.
With three labels, the Stefan model fitted to the target labels by least squares had higher error than the recalibrated model in the Lena Delta (2.9~cm), Canada (5.4~cm) and E Russia (3.7~cm), and lower error in W Russia (5.0~cm) and on the Tibetan Plateau (81.8~cm), where the source coefficient was far from the regional value (descriptive comparison; Fig.~3d).
[XB: SI, R3 지시 문장 | 결과 열람 뒤 설계, 실행 전 등록(비맹검 부분 포함) | 2.2 사전 고정 해석 문장 XB-1, XB-2(라벨 10개 행)]

\subsection{Method selection with tens to hundreds of labels}\label{sec:results-selection}

Selecting the method by five-fold block cross-validation within the target labels reduced RMSE relative to a fixed residual recipe ($\lambda = 0.25$) by 1.3~cm with 160 labels (95\% CI 0.7 to 1.7; block-equal 0.8 to 1.8; Lena Delta and Canada; unchanged under common resampling) and was non-inferior at 40 and 160 labels (margin 0.5~cm), with both results arising in Canada; this rule was designed after the label-grid results had been viewed (Fig.~4c; Fig.~7d).
With 40 labels, the reduction of 0.87~cm (95\% CI 0.27 to 1.20) depends on the split-independence assumption.
In Canada the rule reduced RMSE by 2.1 and 2.7~cm with 40 and 160 labels; in the Lena Delta no difference was established.
With ten labels in the four main regions, non-inferiority was not established (−0.25~cm; 95\% CI −1.11 to 0.73), and the rule had 0.70~cm higher error in the Lena Delta (95\% CI 0.56 to 1.03); with all labels, non-inferiority depends on the split-independence assumption.
The rule did not increase the number of targets with lower error than the recalibrated Stefan model (3 against 4 for the fixed recipe with 160 labels, 2 against 8 with all labels; Supplementary Table S8).

In Canada, the rule chose residual ML with $\lambda = 1.0$ in 22 of 25 draws with 40 labels and in 24 of 25 with 160 labels, and no difference from the source-coefficient Stefan model was established.
With Alaska held out, the rule had higher error than the source-coefficient model with 10, 40 and 160 labels (0.91, 0.98 and 2.1~cm; descriptive curve contrasts).

At 40 and 160 labels, residual ML had 1.9 to 2.4~cm higher error than physics-input ML in the Lena Delta and Canada (the 2.4~cm contrast depends on the split-independence assumption) (Fig.~3c).
This difference arose in Canada (4.5 to 5.1~cm); no difference was established in the Lena Delta, and with Alaska held out residual ML had lower error.
At these label counts no difference was established between residual ML and a stronger recalibrated baseline (the soil-property Stefan model recalibrated with $\kappa = 10$).

Hyperparameter selection for neural networks by leave-one-source-region-out cross-validation gave learner-specific results (Supplementary Table S4): where both intervals agreed, it increased the error of direct ML for the MLP and the multi-head MLP (1.3 to 3.0~cm), no learner met the registered condition for a learner-general statement, and this selection was not compared with selection within the target labels.
[XC: R4 | 결과 열람 뒤 설계, 재사용 지역의 재검정(비맹검 부분 포함) | 2.3 사전 고정 해석 문장 XC-1, XC-2(+XC-2s), 뒤에 S-XC4 (b) 지역 손해 문장]
[XB: SI, R4 지시 문장 | 결과 열람 뒤 설계, 실행 전 등록(비맹검 부분 포함) | 2.2 사전 고정 해석 문장 XB-1, XB-2(라벨 40, 160 행, 레나·캐나다 2지역)]

% results_b.tex : Results, subsections R5 to R8 (body text only, no preamble).
% Spec: paper/manuscript/MANUSCRIPT_SPEC.md sections 4.5 to 4.8 and 11.
% Order follows the binding style guide 4.3: R5 joins the label-rich gains and
% the grid-level gain source, R6 placement, R7 independent regions, R8 intervals.
% Numbers: copied from paper/claims/{C4,C6,C8,C1,C3,C2,SI_learners_new_regions,
% SI_uncertainty}/README.md and rounded with style guide 4.5 (|delta| < 1 cm two
% decimals, other cm values one decimal, percentages integer, CI ends with the
% digits of the printed point estimate). Replace with numbers.tex macros when
% build_numbers.py exists (QA item M-09).
% Sign: negative values use U+2212. A xelatex test build rendered them; a local pdflatex
% test with T1 bitmap fonts compiled without error but dropped the glyph, so check the
% minus signs in any pdflatex build.
% Placeholders [X?: ...] are the registered strings of spec section 11 and must
% be replaced by the pre-fixed interpretation sentences only. They contain Korean
% text, so a build that keeps them needs xeCJK (spec build test 4).
% Citation keys: paper/manuscript/en/refs_results_b.bib.

\subsection{Gains within label-rich regions}
\label{sec:results-label-rich}

Within Alaska, residual ML on the recalibrated Stefan anchor reduced RMSE relative to the Stefan model recalibrated on the same labels by 0.54~cm with 1000 labels (95\% CI 0.32 to 0.69; block-equal 0.40 to 0.80; 14.4 to 13.9~cm; 25 splits; unchanged under common resampling), and with all labels about 99\% of the reduction in squared error lay between ERA5-Land climate grid cells\cite{munozsabater2021_era5_land} (Fig.~7a,c).
The test used no other source region and a residual weight $\lambda$ selected by cross-validation, and it was designed after earlier results were viewed and registered before running.
A first test (5 splits, 1000 or more labels) had established no difference against a physics calibration selected by cross-validation; the second test changed the comparator, the splits (5 to 25), the recipes and the label range (adding 200 and 500 labels) (Methods).

Above the error floor of 11.3~cm, the 1000-label reduction removed 19\% of the reducible squared error (79.7 to 64.4~cm$^2$).
Residual ML also had lower error with 500 and with all labels (−0.39 and −0.52~cm), but these two results depend on the split-independence assumption, and against the source-coefficient Stefan model no difference was established with 1000 or all labels (Supplementary Table S7).
In the Lena Delta (about 1480 labels) and Canada (about 380), no difference was established for the same contrast at any tested label count, and nine synthetic aperture radar covariates left the Alaskan error of residual ML equivalent within ±0.5~cm at 500 to 5000 labels (five splits; no difference established with all labels).

Grouping scoring cells with the same ERA5-Land thawing degree-days within a 0.5° block, a proxy for climate grid cells, split the error into between-grid and within-grid parts (Methods; designed after earlier results were viewed and registered before running).
With all labels, residual ML reduced the between-grid RMSE relative to the recalibrated model by 1.0~cm (95\% CI 0.6 to 1.3; block-equal 0.6 to 1.2), and the within-grid RMSE was equivalent within ±0.5~cm (−0.01~cm; 95\% CI −0.03 to 0.04; block-equal 0.01 to 0.06), with the same values to two decimal places under an air-temperature and a soil-temperature grid definition (Fig.~7c; RMSE contrasts of the two parts do not add).
Within-grid variation made up 72\% of the squared error of the recalibrated Stefan model, but residual ML explained less than 1\% of it in Alaska and less than 8\% in every stored configuration across the three regions.
Attributing the gain to grid-scale bias correction is a post hoc description (error change by scoring block in Fig.~7b).

In the registered within-grid hypothesis, the within-grid predictions of residual ML increased the error in the stratified regional mean with all labels (+0.09~cm; 95\% CI 0.02 to 0.40; block-equal 0.25 to 0.50) and by 0.27 to 0.42~cm with 500 to all labels under the soil-temperature definition, all statistically distinguishable but below 0.5~cm in size; the error was higher in Canada (+0.23~cm; 95\% CI 0.16 to 0.34; block-equal 0.12 to 0.31).
In transfer, there was no evidence that within-grid predictions matched the observed variation.
Extrapolation to blocks warmer than the training range could not be evaluated, because the Canadian extrapolation set had three scoring blocks in every split (Supplementary Table S10).
[XE: R5 SI 지시 문장 | 결과 열람 뒤 설계 | 2.5 사전 고정 해석 문장(6갈래)과 공통 문장]
[XI: R5 SI 지시 문장 | 결과 열람 뒤 설계 | 2.9 다섯 갈래, 공간 대용과 외삽 폭]
[XB: SI, R5 지시 문장 | 결과 열람 뒤 설계, 실행 전 등록(비맹검 부분 포함) | 2.2 사전 고정 해석 문장 XB-3, XB-4]

\subsection{Placement of new labels}
\label{sec:results-placement}

In Canada, spreading labels over 0.5° blocks or covariate space reduced RMSE relative to random cell sampling by 2.4~cm with 40 transfer labels (95\% CI 0.5 to 3.4; block-equal 0.8 to 2.5; covariate max-min distance sampling) and by 2.5 to 3.2~cm with 50 to 100 within-region labels, whereas in the Lena Delta block stratification increased RMSE by 1.7~cm with 100 labels (95\% CI 0.9 to 2.4; block-equal 0.1 to 1.6) (Fig.~5).
All contrasts used residual ML ($\lambda = 0.25$) on the same splits and were designed after earlier results were viewed and registered before running.
The Canadian reduction was 1.3 to 1.4~cm with 200 labels, and in the Lena Delta block stratification gave positive differences at every label count from 20 to 500 (+1.5 to +2.0~cm; cell-weighted intervals above zero).

In Alaska as a whole, no difference was established for either strategy at 50 to 200 labels: block-equal differences were −1.1 to −1.7~cm with intervals below zero, but cell-weighted point estimates were +0.28 to +0.74~cm.
In two Alaskan sub-regions, which are not independent of Alaska, the 12 combinations of sub-region, strategy and label count gave lower error in 6 (0.32 to 0.77~cm), equivalent error in 1 and no established difference in 5.
With ten transfer labels, neither strategy had higher error in any of five regions, which is not a non-inferiority result: block stratification in the Lena Delta had a cell-weighted difference of +1.9~cm (95\% CI 0.3 to 3.1), and covariate spreading in E Russia a block-equal difference of +0.36~cm (0.08 to 0.62) (Fig.~5f).

Sequential selection of the cells with the largest prediction variance increased RMSE relative to random sampling in all nine combinations of Alaska, the Lena Delta and one Alaskan sub-region at 50 to 200 labels (+0.12 to +2.8~cm), and reduced it in Canada at 100 and 200 labels (1.5 and 0.68~cm), less than spreading did at the same label counts.
The strategy with the lowest Alaskan RMSE at 100 labels, covariate cluster stratification, increased RMSE in the Lena Delta at the same label count (+0.33~cm; 95\% CI 0.04 to 0.77; block-equal 0.36 to 1.00; statistically distinguishable but below 0.5~cm in size).
The internally pre-registered contrast of block stratification against random cell sampling, averaged over regions, was undetermined (10 labels: −0.07~cm, 95\% CI −0.86 to 0.89; 40 labels: −0.24~cm, −1.19 to 0.99; two-stage resampling; partly non-blind); it uses the same label sets as the transfer contrasts, and its 40-label mean combines Canada (−2.6~cm) and the Lena Delta (+2.1~cm) (Supplementary Table S9).
We also tested a placement procedure fixed before evaluation (Methods) within the sequential workflow (Fig.~7e).
[XC: R6 | 결과 열람 뒤 설계, 재사용 지역의 재검정(비맹검 부분 포함) | 2.3 사전 고정 해석 문장 XC-5와 공통 문장]
[XD: R6 SI 지시 문장 | 결과 열람 뒤 설계, 탐색(SI), XD-2 사후 설계(지역 표적) | 2.4 사전 고정 해석 문장(SI) XD-1, XD-3, 공통 한계 문장]

\subsection{Independent regions}
\label{sec:results-independent}
% [DECISION: 러시아 중부의 영문 표기] (spec 4.7, 13.1 item 9)
% [DECISION: 지역 이름 표기] (W Russia, E Russia, Lena Delta; spec 1.4, 13.1 item 7)

In Central Russia, a licence-verified new region with a shallow active layer (57 labels), direct ML without target labels increased RMSE relative to the source-coefficient Stefan model by 2.7~cm (95\% CI 2.1 to 4.1; block-equal 3.2 to 6.6; unchanged under common resampling), and in the five-region shallow pool that includes it, residual ML with all labels reduced RMSE relative to the same model by 3.2~cm (95\% CI 2.2 to 4.2; block-equal 1.5 to 3.7) (Table~1).
[MISSING: 5지역 풀 대비(라벨 전량 잔차 ML − 원천 계수 Stefan, 라벨 10개 잔차 ML − 재보정 Stefan)의 지역 중앙값, 오차 감소·증가 지역 수, 러시아 서부 제외 값. 새 지역 판정표(약관 확인분 판)의 지역 행에서 수치 생성 스크립트로 만든다]
These results carry a cross-environment flag, because the pools combine regions fitted on two computing platforms (Methods).
No region in the four-region or the five-region pool had lower error with direct ML than with the source-coefficient model at 0 to 40 labels.
With ten labels, residual ML and the recalibrated Stefan model were equivalent within ±0.5~cm in the five-region pool (−0.13~cm; 95\% CI −0.42 to 0.01; block-equal −0.41 to 0.01), so the four-region result at this label count did not hold there.
In Central Russia alone, no difference was established for residual ML with all labels (cell-weighted −5.4~cm, 95\% CI −9.0 to −1.6; block-equal −2.5~cm, −6.8 to 1.5).

The Tibetan Plateau, with a deep active layer, is reported separately and not used as independent-region evidence (scored cells outside the source elevation range; source-coefficient RMSE 244.5~cm; non-blind contrasts).
There, direct ML without labels reduced RMSE by 60.0~cm, and residual ML with ten labels had a 20.9~cm lower RMSE than the shrunk recalibration but was not better than a least-squares Stefan fit to the same labels (+19.1~cm; 95\% CI 5.7 to 30.1; block-equal −11.8 to 16.8).
The North Atlantic region, with partly unconfirmed data licences, appears only in the full-data edition of the Supplementary Information (Supplementary Table S11).
[XF: R7 | 결과 열람 뒤 설계, 새 지역은 맹검(독립 지역 확인), NAtlantic v5 는 민감도 | 2.6 사전 고정 해석 문장]
[XC: R7 | 실행 전 등록 확인 시험(맹검) | 2.3 XC-F3 갈래]
[XJ: R7 SI 한 줄 | 결과 열람 뒤 설계, 서술 | 2.10 사전 고정 해석 문장]

\subsection{Prediction intervals}
\label{sec:results-intervals}

Without target labels, hierarchical conformal intervals\cite{dunn2023_hierarchical_conformal} calibrated with regions as groups covered on average 0.86 of held-out observations (95\% CI 0.80 to 0.92; nominal 0.90; four regions), from 0.75 in W Russia to 0.93 in E Russia, and with four calibration regions they carry no finite-sample guarantee (internally pre-registered under an earlier protocol; Supplementary Table S12).
Quantile intervals calibrated in Alaska by conformalized quantile regression\cite{romano2019_conformalized_quantile_regression} covered 0.32 to 0.73 in the four held-out regions, and generative quantile intervals calibrated on exchangeable source cells covered 0.76 (normalizing flow) and 0.65 (flow matching).
Across five ways of scaling the width of the hierarchical interval, coverage was 0.86 to 0.95 and mean width 81.3 to 132.4~cm (Fig.~6d).
For the interval score\cite{gneiting2007_proper_scoring_rules}, no difference was established between cell-specific flow scaling and a placebo with widths permuted within regions (+3.3~cm; 95\% CI −4.2 to 8.2), or between any scaling and constant width.
With ten labels, hierarchical predictive distributions and calibrated constant-width intervals differed in interval score by +9.5~cm (95\% CI −0.8 to 22.1; Holm-adjusted \textit{P} = 0.277), so no difference was established.
With 40 and 160 labels, residual-ML intervals covered 0.84 to 0.90 in the Lena Delta, Canada and Alaska (Fig.~6e).


% =====================================================================
% Discussion (English manuscript). Body text only, no preamble.
% Written 2026-10-04 from paper/manuscript/MANUSCRIPT_SPEC.md section 5
% (paragraphs D1-D6, no subheadings, target 850 words).
% Bibliography keys: paper/manuscript/en/refs_discussion.bib.
%
% Number provenance (to be replaced by numbers.tex macros once
% build_numbers.py exists; rounding follows style guide 4.5):
%  D1  2.5 cm (95% CI 1.9 to 3.0), 1 of 4 regions (Holm P 0.001 in Results)
%        = AB4, C2 README 2.1 A6; C3 README E36a
%      non-inferior at 40 and 160 labels, gain in Canada
%        = WF4-a, C5 README 2.1 E2, E3; 2.2 E6, E7
%      point estimates above source Stefan, 3 to 160 labels, pool
%        = FIGURE_SPEC_v3 9.3 content risk; C5 README 2.4 E21-E25
%      about 99% between-grid share = C8 README 2.2 B4 (0.9% within)
%  D2  78-94% Alaskan source cells = D README E35
%      85% of four-region mean reduction = C2 README 2.3 ratio table
%      13.5 of 14.0 cm, W Russia = C2 README RW5, RW1
%      E 1.60 vs 1.59 = C2 README CA6
%      +2.3 cm Canada, ten labels = C2 README CA1 (AB4 region row)
%      residual weight 1.0 in 22 of 25 runs (40 labels) = C5 README E29, N6
%      rho 0.38 (0.17 to 0.55), 28 targets = C2 README A3 (WF4-c)
%  D3  8%, 39%, best of 56 settings (earlier analysis, 14.46 -> 13.33 cm)
%        = C4 README 2.6 E8.contest.*, D2.contest.* (7.80%, 38.60%), 5.1-6
%      4%, 19% (14.41 -> 13.87 cm) = C4 README E1.R1.n1000, D1.wf6.*
%        (3.75%, 19.17%)
%      structure comparison = C3 README 1.3 items 1-2, 1.4
%      13,606 rows, 343 locations = D README E10
%  D4  placebo contrasts, 0.48 cm linear thaw index = C1 README 2.3 (L15)
%      0.86 cm label unit (median, Canada, ten labels) = D README E46
%      single-year labels = D README E47
%      3.6e-14 cm, 5.6 cm = SI_learners README 1.3, 2.6 (R-C14, R-gate-iii)
%  D5  10.3 cm (7.3 to 12.8) = C7 README A1
%      0.39 to 0.54 cm = C4 README 5.1 item 5
%      <1% within-grid explained in Alaska = C8 README 2.3 (C2 max 0.7%)
%      1990-2024 vs 2015-2020 = D README 5 item 5
%  D6  0.70 cm Lena Delta, ten labels = C5 README 2.2 E12
%      72% within-grid share of recalibrated Stefan SSE = C8 README B1
% =====================================================================

\section{Discussion}\label{sec:discussion}

The results define a label-count workflow for a new permafrost region (Fig.~7d). Without target labels, the reference is the source-coefficient (or year-matched) Stefan model, and ML enters only through physics pseudo-labels or a low-weight residual. With three to ten labels, the coefficient is recalibrated and a low-weight residual added; ten-label recalibration lowered the four-region mean RMSE of the Stefan model by 2.5~cm (95\% CI 1.9 to 3.0; lower error in one of four regions), and no further difference was established for the residual. From 40 labels, the method is selected by cross-validation within the target labels, which was non-inferior to a fixed residual recipe at 40 and 160 labels in the Lena Delta and Canada (margin 0.5~cm), with the gain arising in Canada. In this pool, the methods recommended from 3 to 160 labels had higher point-estimate error than the source-coefficient model [DECISION: Fig 7d 경로의 y 기준과 풀, 그림 명세 D-13]. Within regions with thousands of labels, residual ML follows; in Alaska about 99\% of its squared-error reduction lay between climate grid cells (Fig.~7a--c). Placement enters only as a pre-specified procedure, as spreading labels lowered error in Canada but not in the Lena Delta (Fig.~5).

The gains came from correcting the level of the Stefan coefficient $E$ and bias at the climate-grid scale. $E$ lumps soil thermal conductivity, the n-factor and the latent heat of soil water and ice~\cite{Riseborough2008}, and regional Stefan maps have therefore used land-cover-specific coefficients~\cite{Nelson1997,ShiklomanovNelson2002}; here it also absorbs thawing-index bias from grid and year mismatch and the label definition. With all labels, recalibration made up 85\% of the four-region mean reduction by residual ML relative to the source coefficient, a mean dominated by W Russia (recalibration 13.5 of 14.0~cm; post hoc arithmetic). The regional coefficient of Canada was the closest of the main regions to the source value but varied between label and scoring blocks; ten-label recalibration increased error there by 2.3~cm, and the selection rule mostly chose a residual weight of 1.0, which we read as correcting label-location heterogeneity. With all labels, the gain of label-based correction (recalibration plus residual) was rank-correlated with the magnitude of the ten-label bias (Spearman $\rho = 0.38$; 95\% CI 0.17 to 0.55; 28 targets). [XA: Discussion D2 | 사후 분석(재현(비맹검)) | 2.1 해석 조각 XA-3]

The zero-label results agree in direction with an Alaskan random-split comparison in which a random forest generalized less well than the Stefan model~\cite{Gautam2025}. An earlier within-Alaska analysis [DECISION: 대회 보고서 인용 여부] found 8\% lower RMSE and 39\% lower reducible squared error with a physics residual (best of 56 settings), and the stricter test here found 4\% and 19\%. In our tests, feeding physics outputs to ML~\cite{Pilyugina2025,WangG2025} gave higher four-region mean error than the residual structure with ten or fewer labels, and lower error in Canada with 40 and 160 labels. Random cross-validation is optimistic for spatial mapping~\cite{Ploton2020,MeyerPebesma2022}, but only probability samples give design-unbiased map accuracy~\cite{Wadoux2021}; our labels are a clustered non-probability sample (13,606 Alaskan rows at 343 locations), so our scores are prediction errors on held-out blocks, and kNNDM~\cite{Linnenbrink2024} only mimics map-to-label distances. Label-count designs exist outside ALT, for lake temperature against calibrated process models~\cite{Read2019,Jia2021} and for borehole temperature from 1 to 40 local observations in new regions~\cite{OMalley2026}.

Alternative explanations are unlikely. Stefan pseudo-labels had lower error than shuffled and constant pseudo-labels (Fig.~3b), whereas a linear thawing-index placebo was 0.48~cm worse without labels and equivalent with ten labels, so the gain reflects cell-level thawing-index information more than the square-root form. Recalibration gains held with single-year labels in W and E Russia and were 0.86~cm smaller with point labels than with 1~km location means in Canada (median, ten labels), which sets the label unit of the workflow. Physics and CatBoost results of the workflow methods agreed between computing environments (maximum difference $3.6\times10^{-14}$~cm), whereas a neural network differed by up to 5.6~cm, so the workflow is built from these deterministic methods.

The workflow does not rank sites for measurement or predict variation within climate grid cells, so 1~km maps display a climate-grid-scale correction; its conclusions are restricted to source data that include ground-penetrating radar labels, and climates warmer than the training range could not be evaluated. The Alaskan dominance of the source data (78--94\% of source cells) matters most at small label counts, because the source coefficient carries the weight of ten labels in the recalibrated coefficient. Confirmatory evidence comes from four re-used main regions and one new shallow region; the unblinded Tibetan Plateau is the only deep-regime region and is reported separately, so results are not generalized beyond these regions. Values depend on the test design: the RMSE of direct ML differed by 10.3~cm between random cell splits and region holdout (five-region mean; 95\% CI 7.3 to 12.8), so randomly split studies are not directly comparable. Soil moisture, organic-layer thickness, snow duration and microtopography within grid cells are nearly absent from the covariates, so the within-grid share explained in Alaska (below 1\%) describes these inputs, not a physical limit. Method selection, placement, the label-rich tests (gains of 0.39 to 0.54~cm), the bias correlation and the gain decomposition were designed after earlier results were viewed and are weaker evidence than the internally pre-registered contrasts.

Each step carries a condition: labels are 1~km location means, recalibration rests on the ten-label contrast, and selection by cross-validation within the target labels depends on the number of 0.5$^{\circ}$ blocks covered by the labels and increased error relative to the fixed recipe in the Lena Delta with ten labels (0.70~cm). [XC: Discussion D6 | 결과 열람 뒤 설계, 재사용 지역의 재검정(비맹검 부분 포함) | 2.3 XC-1, XC-2] The data needed next are independent regions not used in development [XF: Discussion D6 | 결과 열람 뒤 설계, 새 지역은 맹검(독립 지역 확인) | 2.6 해석 문장], existing ALT maps scored on the same blocks~\cite{Ran2022a,Wei2026} [XG: Discussion D6 | 결과 열람 뒤 설계, 등록 이탈(WRAPUP 10) | 2.7 공통 문장], and above all covariates that vary within climate grid cells, because 72\% of the recalibrated Stefan squared error in Alaska lay within those cells while the gain of ML lay between them. [XE: Discussion D6 | 결과 열람 뒤 설계 | 2.5 해석 문장]


% Methods; its last subsection is Code availability (front.tex block 2 is the same subsection and is not used)
% methods.tex : Methods (English manuscript, body text only, no preamble).
% 작성 2026-10-04. 명세: paper/manuscript/MANUSCRIPT_SPEC.md 6절(M1-M15, 약 4240단어 목표), 1.4절(방법 이름),
% 1.5절(서술 방향), 8절(Supplementary Methods 번호), 11절(자리표시 등록부), 12절(QA).
% 구속 지침: design/journal_grade_style_guide.md 4절. 인용 키: paper/manuscript/en/refs_methods.bib
% (refs_discussion.bib 와 같은 FirstAuthorYear 규칙).
% 방법 약호(P0, P1, R1, D0, D1, W)는 쓰지 않는다. 실험·가설 번호(LG, LGX, WF, XA-XJ, L1-L43, AB1-AB10)는
% 등록 소절(Internal pre-registration and deviations, Additional registered analyses)에만 둔다.
% 코드 서체(\texttt)는 Code availability 안에만 쓴다. 세부는 명세 8절의 Supplementary Methods 1-9 로 보냈다.
%
% 수치 출처(원고 조립 때 shared/numbers.tex 매크로로 바꾼다. 자릿수는 지침 4.5):
%   17,572 / 17,467 / 68 .............. claims/D_data_and_design/README.md E03, E04 (68 = v3 출처 표지 변경 행 = F4_calm_temp)
%   20 (15 / 4 / 1) ................... D README E08, E50
%   30 / 13,606 / 343 / 74 ............ D README E14, E10
%   39.7 / 15.1 / 8.7 rows per 1 km ... D README E21-E23 (39.67, 15.11, 8.72)
%   57 rows, 13 blocks / 132, 34 / 41 . D README E17-E19
%   78-94 % / 1.59-1.62 (E0, main targets; Alaska(x) 1.50 has a Lena-dominated source) ... D README E35, E37
%   14-17 cells in A .................. D README E29, E30
%   21.7-43.0 cm (P0 RMSE, cell) ...... D README E39 (21.69-42.98); '이번 계산' 1.2-2.3 % -> '약 1-2 %'(명세 6.9)
%   13,333 cells (error floor) ........ D README E45
%   13-74 blocks ...................... D README E10-E14, E17 (Russia C expanded 13; 2026-10-04 check)
%   27 targets, 5 splits, 5 draws, 2 seeds, 12 methods ... D README E24, E25
%   adjusted Rand 0.999 / 0.997 / 0.832  claims/C8_gain_source/README.md 5.1 의 1
%   3.6e-14 cm (34,360 keys); 11,814 keys; 43 of 288, 5.61 cm; 0.46 cm ... claims/SI_learners_new_regions/README.md 2.6
%   commit hashes and times ........... git log (committer time, KST); docs/MANUSCRIPT_DRAFT_SUPPORT_2026-09-30.md M4.2
%   selection rule (7 candidates) ..... claims/C5_method_selection/README.md 1.4
%   Algorithm P ....................... docs/EXPERIMENT_PLAN_FINAL_BATCH_2026-10-04.md 2.3, 2.4; docs/QA_FOLLOWUP_2026-10-04.md 2.4
%   placement strategies .............. claims/C6_label_placement/README.md 3.1
%   within-region retest (four changes)  claims/C4_sufficient_labels/README.md 1.2 (c)
%   XA-XJ designs ..................... EXPERIMENT_PLAN_FINAL_BATCH 0.1, 1, 2.1-2.10, 6
% 자리표시: [DECISION: ...] 6개, [미확인: ...] 4개, [XD: Methods M7 | ...] 1개(명세 11절 문자열 그대로).
% 자리표시에 한글과 γ 가 있으므로 초안 컴파일은 xelatex(+xeCJK 또는 kotex)로 한다. 제출 전 점검에서 0건이어야 한다.

\section{Methods}\label{sec:methods}

\subsection{Study regions and labels}\label{sec:m-labels}

Active-layer thickness (ALT) labels of the main regions came from the GTN-P/CALM records\cite{Streletskiy2025}, the ALLena compilation for the Lena River Delta region\cite{Veremeeva2025} and version 2 of the ABoVE data set for Alaska and Canada\cite{Moore2025}. Records were aggregated to cells so that repeated measurements at one location do not enter as independent labels. A cell is a coordinate rounded to four decimal places (ABoVE and ALLena) or a CALM site, and its label is the mean of its observations (Supplementary Methods 1). The cell table has 17,572 rows, of which 17,467 carry a direct label from mechanical probing or ground-penetrating radar (GPR). Sixty-eight CALM cells whose metadata name ground temperature or thaw tubes are not used as labels. Twenty temperature-derived cells (15 in Alaska, 4 in Canada and 1 on the Tibetan Plateau) were found after the table had been frozen and remain in the direct class with a flag. Label years span 1990 to 2024, so predictions are static and refer to the 2015--2020 covariate climatology.

Macro-regions were defined geographically before the analysis, namely Alaska, the Lena Delta, Canada, W Russia, E Russia, Central Russia and Greenland [DECISION: 지역 이름 표기] [DECISION: 러시아 중부의 영문 표기]. Among Alaska and the four main regions, label rows range from 30 in E Russia to 13,606 in Alaska (Table~1). Ten sub-regions (Supplementary Table S1) were defined by $k$-means clustering of 0.5$^{\circ}$ block centroids without labels and were fixed before the main run.

The label count $n$ is a row count, and the meaning of a row differs between regions. A row is the multi-year mean of a CALM site in W Russia and E Russia, usually a single-visit observation in the Lena Delta, mostly a point observation in Alaska and Canada, and a mean over about 1~km in the additional regions. Alaska, the Lena Delta and Canada have 39.7, 15.1 and 8.7 rows per 1~km location, and the 13,606 Alaskan rows lie at 343 locations in 74 blocks.

Additional regions were assembled from public archives under rules fixed before any fit (Supplementary Methods 2). A label had to be a direct end-of-season thaw depth from 1990 or later, at least 100~km from existing labelled regions, and was aggregated to a cell of about 1~km. A region was eligible for confidence intervals if the scoring blocks of its valid splits formed a union of at least eight blocks. Sources without a confirmed redistribution licence were removed in a licence-verified edition that defines the main verdicts (Supplementary Table S17). In this edition, Central Russia (57 rows, 13 blocks, shallow regime) and the Tibetan Plateau (132 rows, 34 blocks, deep regime) are eligible. A North Atlantic group (41 rows) is eligible only in the full edition and is reported in the Supplementary Information.

\subsection{Covariates}\label{sec:m-covariates}

All learning methods use 25 covariates that are available in every region. The six terrain covariates are elevation, slope, the sine and cosine of aspect, and the topographic position index and roughness in windows of $33 \times 33$ pixels, from the Copernicus DEM GLO-30 [미확인: Copernicus DEM GLO-30 인용 형식]. The eight climate covariates are mean annual air temperature, thawing degree-days (TDD), freezing degree-days, $\sqrt{\mathrm{TDD}}$, warmest- and coldest-month air temperature, level-1 soil temperature and snow water equivalent, from ERA5-Land monthly means for 2015--2020\cite{munozsabater2021_era5_land}. The nine soil covariates are clay, sand and silt content, bulk density, coarse fragments and pH at 5--15~cm and soil organic carbon at three depths, extracted from the 250~m SoilGrids 2.0 layers\cite{Poggio2021} resampled to about 5~km. The two model-based covariates are the 1997--2021 mean of the ESA CCI Permafrost ALT product, version 4.0, and its validity flag\cite{Westermann2024}. This product is model output forced partly by the same reanalysis and is not treated as an observation.

TDD and the Stefan solution\cite{Stefan1891} are
\begin{equation}
\mathrm{TDD} = \sum_{m} \max(\bar{T}_m, 0)\, d_m, \qquad \mathrm{ALT} \approx E\, s, \qquad s = \sqrt{\mathrm{TDD}},
\label{eq:stefan}
\end{equation}
where $\bar{T}_m$ is the monthly mean 2~m air temperature of month $m$, $d_m$ is its number of days and $E$ is in cm\,($^{\circ}$C\,d)$^{-1/2}$. A cell is scored only if its label, its CCI value and its TDD from level-1 soil temperature are valid, and all methods share the scoring cells. Quantities derived from labels, such as the count and years of observations, are excluded from the inputs, and unit tests enforce this. The information resolution of the inputs is about 0.1$^{\circ}$ for climate and about 5~km for soil.

\subsection{Stefan baselines and coefficient recalibration}\label{sec:m-stefan}

The coefficient $E$ absorbs soil thermal properties, water content and surface conditions, and regional applications calibrate it against local measurements\cite{Nelson1997}. The source set $S$ of a target contains all labelled cells outside the target region and a 100~km buffer around it, and the source coefficient is the least-squares slope through the origin,
\begin{equation}
E_0 = \frac{\sum_{i \in S} s_i\, y_i}{\sum_{i \in S} s_i^2},
\label{eq:e0}
\end{equation}
where $y_i$ is the observed ALT. The source-coefficient Stefan model predicts $E_0 s$. Alaskan cells form 78--94\% of the source sets of the main targets, so $E_0$ varies little between these targets (1.59--1.62~cm\,($^{\circ}$C\,d)$^{-1/2}$; Table~1). With $E_{\mathrm{ls}}$ the same estimator computed from the $n$ target labels, the recalibrated Stefan model predicts $E_n s$ with
\begin{equation}
E_n = \frac{n\, E_{\mathrm{ls}} + \kappa\, E_0}{n + \kappa}, \qquad \kappa = 10.
\label{eq:en}
\end{equation}
The source coefficient thus carries the weight of ten labels, because recalibration with few labels can increase error\cite{Steyerberg2004}. The value of $\kappa$ was fixed from earlier experiments of this project and is not blind, and $\kappa = 3$ and $\kappa = 30$ are sensitivity cases. The local least-squares Stefan model uses $E_{\mathrm{ls}}$ without shrinkage. A gain of machine learning (ML) is stated only against the recalibrated model fitted to the same labels, because the gain over the source coefficient includes the gain of recalibration.

A registered comparison of physics models and ALT products added two stronger baselines. The year-matched Stefan model applies $E_0$ to $\sqrt{\mathrm{TDD}}$ averaged over the observation years of each label within 2010--2024 (2010--2014 for earlier labels) and is reported at zero and at all labels. At 40 and 160 labels, the output of a soil-property Stefan model (Supplementary Methods 3), recalibrated with equations~(\ref{eq:e0}) and~(\ref{eq:en}), is reported as the stronger recalibrated baseline.

\subsection{Machine-learning methods and combination structures}\label{sec:m-ml}

In the transfer design, training rows are the source cells and the $n$ target labels with equal weight. Anchored methods have the form
\begin{equation}
\hat{y}(x) = a(x) + \lambda\, g(x),
\label{eq:anchor}
\end{equation}
where $a$ is a physics anchor and $g$ is a learner trained on residuals from it. Direct ML regresses ALT on the covariates, which include $\sqrt{\mathrm{TDD}}$ and the CCI value. Physics pseudo-label augmentation adds rows with value $E_n s$ at cells of the label half of the target, ten per source row, while drawn cells keep their observed value. Source anchor plus residual ML uses $a = E_0 s$. Recalibrated anchor plus residual ML (residual ML) uses $a = E_n s$ and residual targets $y - E_0 s$ for source rows and $y - E_n s$ for target rows. Augmentation plus anchor plus residual adds zero-valued residual rows at label-half cells. Physics-input ML adds the outputs of a Kudryavtsev-type model [미확인: Kudryavtsev 1974 서지] and of the soil-property Stefan model, or the anchor value, to the inputs of direct ML. The Stefan-CCI anchor is the mean of $E_n s$ and the CCI value where that value is valid. Four further methods of the main experiment are described in Supplementary Methods 3. The residual weight $\lambda \in \{0.25, 0.5, 1.0\}$ is applied after fitting. The reference weight for verdicts, fixed before the main run, is 0.25 for residual structures and 1.0 for direct and physics-input structures, and the source anchor plus residual at 0.25 is called the low-weight residual.

Four placebos tested whether the gain of augmentation comes from the cell-level Stefan values. They replaced the pseudo-labels by the same values shuffled among cells, by their target-level mean, by the mean source ALT, or by a function linear in TDD fitted to the source.

The learner of the method comparisons is CatBoost\cite{Prokhorenkova2018} (200 iterations, learning rate 0.05, depth 3, two seeds). For direct ML without target labels, ten learners were compared. They were this setting, a larger CatBoost, CatBoost tuned by leave-one-source-region-out cross-validation, a random forest, TabPFN v2\cite{Hollmann2025} and TabICL v2\cite{Qu2026} without fine-tuning, and four neural networks with hyperparameters selected by leave-one-source-region-out cross-validation on source data\cite{BergstraBengio2012,CawleyTalbot2010}. The networks are a multilayer perceptron, a multi-head perceptron (not TabM\cite{Gorishniy2025}), a reduced FT-Transformer-type model\cite{Gorishniy2021} and RealMLP\cite{Holzmuller2024} (settings in Supplementary Methods 5). In the figures these methods are labelled Source Stefan, Year-matched Stefan, Recalibrated Stefan, Anchor + residual ML, Source anchor + residual, Physics pseudo-labels, Direct ML, Physics-input ML, Stefan-CCI anchor and, for the rule described below, Target-label CV selection.

\subsection{Transfer design}\label{sec:m-transfer}

To prevent information leakage from nearby labelled cells, the source of a target excludes the target region and a 100~km buffer. Sub-regions were evaluated with the parent region excluded from the source and, as a simulated new area inside a labelled region, with the rest of the parent kept. The 27 targets are the six version-3 regions other than Alaska, Alaska and the ten sub-regions in both modes.

The 0.5$^{\circ}$ blocks of each target were split five times into halves A and B by assigning randomly permuted blocks to the half with fewer cells, so that labelled and scored cells are separated by block boundaries\cite{Roberts2017}. Labels were drawn from A and errors were scored on the masked cells of B. For each $n$ in $\{0, 3, 10, 40, 160, 320, 1000\}$ below the size of A, and for all cells of A, five random draws were made with seeds fixed by target, mode, split, $n$ and draw, so that all methods used the same labels. W Russia and E Russia have 14--17 cells in A, so pooled means at 40 and 160 labels contain only the Lena Delta and Canada.

The four main regions (the Lena Delta, Canada, W Russia and E Russia) form the confirmatory pool. They were used in earlier experiments of this project, so their verdicts are re-tests of re-used regions. Alaska is a reference target, reported separately and in an auxiliary three-region mean with the Lena Delta and Canada. Central Russia and Greenland of the version-3 table are reported as point estimates, and sub-regions, which share cells with their parent, are counted by verdict class. The additional regions enter extended pools of five shallow regions and of six regions including the Tibetan Plateau. All Tibetan scored cells lie above the 99.5th percentile of source-cell elevation, so Tibetan contrasts are not used as evidence for independent regions.

A validation ladder scored direct ML and a Stefan model fitted to the training folds on Alaska and the four main regions under eight schemes of increasing designed separation, from random cells to region holdout (Supplementary Methods 3). An earlier Alaskan analysis added $k$-fold nearest-neighbour distance matching (kNNDM)\cite{Linnenbrink2024}, which matches test-to-training distances to those between a prediction grid and the labels.

\subsection{Within-region design and gain decomposition}\label{sec:m-within}

Within-region tests ask what ML adds where hundreds to thousands of labels exist, and they were designed after the transfer results had been viewed. In Alaska, the Lena Delta and Canada, labels were drawn from half A and were the only training rows. The source coefficient of the transfer design served as prior, $\lambda$ was chosen from $\{0.25, 0.5, 1.0\}$ by block cross-validation within the drawn labels, and augmentation used ten pseudo-labels per label. A first test (5 splits, 1000 or more labels, comparator the best cross-validated physics calibration, nine SAR covariates added in Alaska) was rejected. A second test changed four elements, namely the comparator (the recalibrated Stefan model), the splits (25, of which 24 were valid in the Lena Delta), the recipes (Stefan-CCI anchor and direct ML added) and the label range (200, 500, 1000 and all labels). Both tests are reported.

The gain decomposition asks at which spatial scale ML reduces error. Scoring cells were grouped by 0.5$^{\circ}$ block and air $\sqrt{\mathrm{TDD}}$, so that Stefan predictions and the eight climate covariates are constant within a group. With $N_g$ cells and means $\bar{y}_g$ and $\bar{\hat{y}}_g$ in group $g$,
\begin{equation}
\mathrm{SSE} = \sum_g \sum_{i \in g} \left[(y_i - \bar{y}_g) - (\hat{y}_i - \bar{\hat{y}}_g)\right]^2 + \sum_g N_g \left(\bar{y}_g - \bar{\hat{y}}_g\right)^2,
\label{eq:decomp}
\end{equation}
where the terms are the within-grid and between-grid components. The within-grid RMSE uses groups with at least two cells, and for a Stefan model it equals the within-group standard deviation of the observations. Shares of the reduction are computed from the SSE because RMSE differences of components are not additive. The registration named soil $\sqrt{\mathrm{TDD}}$ for the grouping, whereas the run used the air $\sqrt{\mathrm{TDD}}$ of the Stefan predictions. A soil-grouping sensitivity analysis was therefore registered and run (adjusted Rand index between groupings 0.999 in Alaska, 0.997 in Canada and 0.832 in the Lena Delta).

\subsection{Label placement procedures}\label{sec:m-placement}

Placement strategies select cells of half A without using label values, except for active selection. Besides cell-random sampling, they were block stratification, covariate-cluster stratification, covariate max-min distance sampling, source-extrapolation priority, active selection by the prediction variance of the residual learner and $\sqrt{\mathrm{TDD}}$ priority (Supplementary Methods 8). They were compared within regions at 20 to 500 labels and, for block and max-min sampling, in transfer at 10 and 40 labels. These tests were designed after results had been viewed, whereas a block-spread contrast of the main experiment was registered before results.

Algorithm P combines block allocation with covariate coverage and is the pre-specified placement procedure tested in the workflow analysis. It uses the standardized covariates of candidate cells (no label values), block sizes $N_b$, a nested budget of 10, 40 and 160 labels, an allocation exponent $\gamma$ and a seed.
\begin{enumerate}
\item Order blocks by farthest-first traversal of their covariate centroids $\mu_b$, starting from a seeded block or the block of existing labels.
\item At budget $n_k$, allocate $q_b = n_k N_b^{\gamma} / \sum_{b'} N_{b'}^{\gamma}$ labels to block $b$ by largest remainders, capped at $N_b$.
\item Fill the block with the largest remaining deficit, with its cell nearest to $\mu_b$ if it has no label and otherwise with its cell farthest from all chosen labels.
\item Keep earlier labels when the budget grows.
\end{enumerate}
The candidate settings are $\gamma = 0$, 0.5 and 1, a variant with farthest-first choice of the first cell in a block, block stratification and max-min sampling. One candidate is selected by a mechanical rule from Alaskan targets only, so that the evaluation regions do not inform the choice. Under this rule all candidates were excluded, so Algorithm P is cell-random sampling, and the workflow contrasts of Algorithm P against random placement were not tested. The placement step of the workflow therefore contributed nothing: with ten labels the workflow and its random-placement comparator use the same labels and model, and with 40 and 160 labels the comparison with random placement and the fixed residual recipe compares the selection rule with the fixed recipe on the same labels. Following the registered reporting rule, no map that ranks measurement locations was produced.

\subsection{Method-selection rule and bias diagnosis}\label{sec:m-selection}

The selection rule chooses a method from the target labels alone.
\begin{enumerate}
\item With fewer than ten labels, use the recalibrated Stefan model.
\item Assign the blocks of the labels at random to at most five folds (cell folds if one block). With fewer than two folds, use the recalibrated Stefan model.
\item In each fold, recompute $E_n$ and fit seven candidates, namely the source-coefficient, recalibrated and local least-squares Stefan models, residual ML with $\lambda = 0.25$ and 1.0, augmentation plus anchor plus residual and physics pseudo-labels.
\item Choose the candidate with the lowest cross-validated RMSE (ties to the smaller $\lambda$ and the earlier candidate).
\item Refit it on all labels.
\end{enumerate}
The rule was registered after the label-grid results had been viewed and was applied at 10, 40, 160 and all labels to 30 targets (the 27 transfer targets and three licence-verified additional-region targets). It was compared with a fixed residual recipe with a non-inferiority margin of 0.5~cm.

The bias diagnostic is the absolute mean error of the source-coefficient Stefan model at the ten labels of a draw, averaged over draws and splits. It uses only labels that recalibration also uses, so it is known before any learner is fitted. Its Spearman correlation with the gain of the selection rule over the source-coefficient model at all labels was computed over 28 targets. This gain includes recalibration, so the correlation concerns label-based correction (recalibration plus residual). The part of the gain beyond recalibration was examined in a separate post hoc analysis (Additional registered analyses).

\subsection{Statistical analysis}\label{sec:m-stats}

With $\mathrm{SSE}_b$ and $N_b$ the sum of squared errors and the number of scoring cells in block $b$ of half B, two root-mean-square errors (RMSE, cm) were computed,
\begin{equation}
\mathrm{RMSE}_{\mathrm{cell}} = \left(\frac{\sum_b \mathrm{SSE}_b}{\sum_b N_b}\right)^{1/2}, \qquad \mathrm{RMSE}_{\mathrm{block}} = \frac{1}{|B|}\sum_b \left(\frac{\mathrm{SSE}_b}{N_b}\right)^{1/2},
\label{eq:rmse}
\end{equation}
where $|B|$ is the number of scoring blocks. The cell-weighted value is primary, and the block-equal value limits the influence of densely sampled blocks. A contrast $\Delta$ is the RMSE of the first method minus that of the second, paired by split, draw and seed and averaged over splits, so that negative values mean lower error for the first method. Stratified means weight regions equally.

Confidence intervals (CIs) are 95\% percentile intervals of a bootstrap of the scoring blocks within each split, with the same resampled blocks for both methods. The hypotheses aggregated by the label-grid harness (the main experiment and the additional-region runs) and the prediction-interval experiment used 1000 replicates, and all other contrasts used 10,000. The registered block-spread contrast of the main experiment and the workflow analysis, whose arms use different label sets, used a two-stage bootstrap that also resamples draws; the placement tests report intervals conditional on the label draws. An auxiliary interval resampled the union of a region's blocks once for all splits. Where it changes the verdict, the result is described as depending on the split-independence assumption, and statements without this qualification require both verdicts to agree. A region entered a pooled interval if its valid splits had at least eight scoring blocks in their union.

A contrast is lower error if the upper limits of both the cell-weighted and block-equal CIs are below zero, higher error if both lower limits are above zero, equivalent if neither holds and all four limits lie within $\pm 0.5$~cm (two one-sided tests at 2.5\%\cite{Schuirmann1987}), and undetermined otherwise (``no difference was established''). The margin was fixed before any result. It is about 1--2\% of the cell-weighted RMSE of the source-coefficient Stefan model in the main regions (21.7--43.0~cm) and is a convention, not a measurement error. Margins of 1.0~cm and of 2\% of that RMSE are sensitivity cases. A lower- or higher-error contrast with $|\Delta| < 0.5$~cm is reported as distinguishable with a magnitude below 0.5~cm. Non-inferiority is one-sided and requires both upper limits to be below $+0.5$~cm. The minimum label count $n^{*}$ of a contrast is the smallest $n$ at which the upper limit of the cell-weighted CI is below zero and at least two thirds of splits and 75\% of draws have $\Delta < 0$.

The significance level is $\alpha = 0.05$ two-sided, corresponding to 0.025 one-sided. A $P$ value is the larger of the two-sided bootstrap $P$ values of the two weightings, and unadjusted and Holm-adjusted $P$ values of all registered contrasts are listed in Supplementary Table S2. Ten contrasts fixed before the results are the only ones that may carry verdict words in the abstract. Their $P$ values were adjusted by Holm's procedure\cite{Holm1979} with the family size fixed at ten, and a direction is stated in the abstract only if the adjusted $P$ is below 0.05. Otherwise the abstract states that no difference was established and the main text adds ``significant before correction only''. Results of the within-region, placement and selection tests and of the additional analyses appear in the abstract as numbers only. Within other registered families, Holm-adjusted $P$ values are auxiliary, except in the workflow test and the product comparison, where they set the wording.

Pooled contrasts rest on four to seven regions, so the resampling unit is the 0.5$^{\circ}$ block (13 to 74 per region), and region values and ranges accompany every pooled mean. For statements about regions in general, key contrasts were recomputed over 50 splits in the four main regions and Alaska and summarized by a sign test, a sign-flip permutation test and a Hartung-Knapp random-effects mean\cite{HartungKnapp2001}. Such statements are made only if the Hartung-Knapp interval excludes zero, and with four regions the smallest two-sided sign-test $P$ is 0.125. Rank correlations are Spearman coefficients with intervals from joint resampling of blocks and draws. An error floor, the pooled variance of observations within groups of at least five cells sharing coarse non-terrain covariates (13,333 cells in Alaska), serves only as a denominator of effect sizes.

Analyses designed after earlier results had been viewed are labelled as such, and post hoc descriptions carry no verdict words. Each registered hypothesis carries one of four blinding labels. These are blind, partly unblinded (direction seen under another protocol), replication (unblinded) (direction seen for the same contrast) and results existed at registration (unopened). Prediction intervals were evaluated by empirical coverage and by the interval score $(U - L) + (2/a)(L - y)_{+} + (2/a)(y - U)_{+}$ of a central interval $[L, U]$ with $a = 0.1$\cite{GneitingRaftery2007}. Hierarchical conformal intervals used source regions as groups\cite{Dunn2023} and, with two to four groups, carry no finite-sample guarantee (Supplementary Methods 4). No values are given after $\pm$.

\subsection{Internal pre-registration and deviations}\label{sec:m-registration}

Hypotheses, contrasts, decision rules and permitted sentences were committed to the project's git repository before the results they govern were retrieved (internal pre-registration). The commits were not time-stamped by a third party, and times are committer times in Korea Standard Time (full table in Supplementary Methods 6). The label-grid experiment (LG, L1--L8) was registered in commit 40be64c on 29 September 2026 at 15:45:46. Its extension (LGX), the foundation-model and additional-region runs (LGT, LGD), the interval experiment (LGU), the TabICL and tuned-network runs (LGF) and the reporting rules with the abstract bundle AB1--AB10 (WRAPUP, 316714c, 30 September 2026, 03:02:42) followed, all before the first retrieval of results at 07:20:22 that day.

The workflow experiments were designed after these results had been viewed and were registered before running, WF0--WF5 in 4168331 (1 October 2026), WF6--WF8 in 1b42b4d (1 October 2026), WF9 and WF10 in 8180632 (2 October 2026) and the soil-grouping sensitivity in 7da0064 (2 October 2026). WF0 is a post hoc description, and the priority map planned as WF5 was not produced. The analyses XA to XJ were registered in commit 2678100 (4 October 2026, 14:22:39) before any fit. Only XC-F3 is labelled a confirmatory test registered before running, and XA is a post hoc analysis. The placement of each analysis in the main text or the Supplementary Information was fixed before its results.

Deviations and post hoc changes are as follows.
\begin{enumerate}
\item The gain decomposition (WF9) grouped cells by air rather than the registered soil $\sqrt{\mathrm{TDD}}$, and a soil-grouping sensitivity analysis was added.
\item After the first within-region test (WF1-a) was rejected, the second test (WF6) changed comparator, splits, recipes and label range.
\item The product comparison (XG) was registered as a descriptive table conditional on an interval hypothesis (LGU-B2). After that hypothesis was undetermined, XG was re-registered with six Holm-adjusted contrasts and carries a registration-deviation label.
\item The workflow splits (XC) moved from seeds 6--15 to 201--210 because seeds 1--200 had been used (LGX-N4), and a proposal to use XC verdict words in the abstract was withdrawn.
\item AB4 carries a note that a pre-result sensitivity calculation had displayed the split distribution of the same quantity.
\item Twenty temperature-derived cells found after the table had been frozen remain in the direct class with a flag.
\item Four technical amendments (ff1be02, 30 September 2026, 11:44:30) were committed after the first retrieval but before any result was unpacked or opened, and they change no hypothesis, contrast or verdict rule. [DECISION: 수용 여부, 방법 초안 183행]
\item CatBoost did not pass cross-environment check (i), so extended-pool rows that combine platforms carry a cross-environment flag.
\item No map of measurement priorities was produced, following the registered rule.
\item The climate-extrapolation test (WF10) was not decidable because its extrapolation set had three blocks in Canada.
\item XA resampled eight target families instead of the seven registered, because the North Atlantic target of its 28-target set was not in the registered family list (implementation deviation recorded with the sealed results).
\end{enumerate}

\subsection{Additional registered analyses}\label{sec:m-xbatch}

The ten analyses XA to XJ follow the conventions above, and their placement in the main text or the Supplementary Information was fixed before results (Supplementary Methods 9; results in Supplementary Table S13).

XA (main text, post hoc) splits the gain over the source-coefficient model into the shares of recalibration, of remaining shrinkage and of ML beyond recalibration, and estimates their rank correlations with coefficient error under two definitions, with intervals from joint resampling of target families and blocks. The registered list had seven families; the North Atlantic target, which was not in that list, was resampled as an eighth single-target family, and the 27 targets of the seven registered families were analysed as an unregistered sensitivity with the same result categories. Results are stated only as an interval excluding zero or as not established. XB (Supplementary) learns convex weights of physics models and ALT products from the target labels by cross-fitting, with and without a residual learner. XC (main text) applies Algorithm P, the ten-label diagnostic and the method rule in sequence on new splits of the five shallow regions. Its eight Holm-adjusted hypotheses compare the workflow with random placement and a fixed residual recipe at 40 and 160 labels, test non-inferiority to the recalibrated Stefan model at 10, 40 and 160 labels, compare Algorithm P with random placement under the selection rule, and test the workflow at ten labels in newly obtained regions. The evaluation cells were used in designing the components, so all but the last hypothesis are partly unblinded.

XD (Supplementary) selects the setting of Algorithm P on Alaskan targets, tests alternative coefficient estimators and evaluates a learned placement policy once outside Alaska without $P$ values. XE (Supplementary) adds public inputs at 10--500~m to residual ML within regions and compares within-grid RMSE with the 25-covariate model and the recalibrated Stefan model. XF (Supplementary, blind) applies the eligibility rules to new public data and repeats core contrasts in each eligible region. XG (main text) scores existing ALT maps on the same blocks after masking blocks within 5 or 25~km of product training sites, raw against the source-coefficient model and recalibrated against residual ML. XH (main text) describes, without verdict words, the errors of the recalibrated Stefan model, residual ML and direct ML under random-cell, site, block, kNNDM and region-holdout validation in three regions. XI repeats the climate-extrapolation test on the licence-verified expanded Canadian table, and XJ tests temperature-derived labels as low-weight auxiliary rows (both Supplementary).

\subsection{Computing environments and reproducibility}\label{sec:m-computing}

Physics models and CatBoost run deterministically on CPUs. The label-grid experiment and its extension ran on cloud nodes (Rescale) with 64 CPU cores and four NVIDIA T4 GPUs, and the workflow experiments and additional analyses on cloud CPU nodes of one processor family (AMD EPYC Milan, 96 cores, and Milan-X, 64 cores). The foundation models, tuned networks, interval experiment and additional-region runs used a shared local server with NVIDIA RTX 3090 GPUs. Package versions were pinned (CatBoost 1.2.10 on all platforms; Supplementary Table S16), and each output records code hash, configuration hash, package versions and device.

Physics and CatBoost results agreed between cloud and local runs to within $3.6 \times 10^{-14}$~cm over 34,360 keys, and the two cloud node types gave identical errors over 11,814 common keys. Neural networks did not reproduce across platforms. Of 288 FT-Transformer keys, 43 differed by more than 0.5~cm (maximum 5.61~cm) although refits within a platform were identical, and CatBoost on a ridge anchor differed by up to 0.46~cm. Both sides of a contrast therefore come from one platform, pooled means across platforms carry a cross-environment flag, and the workflow uses only CatBoost and closed-form physics.

\subsection{Maps and colour scales}\label{sec:m-maps}

Maps use a north polar stereographic projection with true scale at 70$^{\circ}$~N and the central meridian stated in each legend, and the Tibetan Plateau appears in a Lambert azimuthal equal-area inset. Land is from Natural Earth (1:50m), permafrost zones from [미확인: 영구동토 구역 자료의 이름과 판], and maps were drawn with Cartopy [미확인: Cartopy 판]. Scientific colour maps\cite{Crameri2018Zenodo,Crameri2020} were used, oslo reversed and restricted to 0.12--0.90 for ALT, broc centred on zero for error differences, bam for residuals and acton reversed and restricted to 0.05--0.95 for interval width.

\subsection{Use of large language models}\label{sec:m-llm}

[DECISION: LLM 사용 기재 문구]

\subsection{Code availability}\label{sec:m-code}

Code for data preparation, experiments, statistics, figures and tests is available at [DECISION: 공개 저장소 주소와 공개 시점] and archived with a version DOI [DECISION: Zenodo DOI]. The main components are the label-grid experiment (\texttt{scripts/3\_deep\_learning/h40\_label\_grid.py}), the validation ladder (\texttt{scripts/2\_evaluation/h41\_validation\_ladder.py}), the extension experiment (\texttt{scripts/3\_deep\_learning/h42\_label\_grid\_ext.py}), the workflow experiments (\texttt{scripts/3\_deep\_learning/h54\_workflow.py}), the abstract-contrast aggregator (\texttt{scripts/2\_evaluation/h39\_scenarios.py}), the post hoc reanalysis (\texttt{scripts/2\_evaluation/wf0\_reanalysis.py}), the additional analyses (\texttt{x\_*.py}), the unit tests (\texttt{tests/}) and the cloud job configurations (\texttt{configs/rescale/}), together with the registration documents and their revision histories. The additional analyses pin catboost 1.2.10, scikit-learn 1.9.1, pandas 2.3.3, scipy 1.17.1 and numpy 2.1.1. The pretrained weights of TabPFN v2 and TabICL v2 are not redistributed and are obtained from their providers.


% Block 3 of front.tex: Data availability, at the end of the main text before the references
\def\frontblock{3}% front.tex : 영문 원고의 제목 쪽, 초록, 핵심어, Code availability, Data availability,
%             Acknowledgements, Author contributions, Additional information(Competing interests)
% 작성 2026-10-04. LaTeX 본문 조각이다(프리앰블 없음). 명세: paper/manuscript/MANUSCRIPT_SPEC.md
%   1.1(순서), 1.3(제목·핵심어), 1.5(서술 방향), 2.1-2.4(초록), 6.15(M15 Code availability),
%   6.16(Data availability), 13.1(사용자 결정 대기).
% 서식: Springer Nature 템플릿 sn-jnl(sn-nature 선택지)의 명령(\title, \author, \affil, \abstract,
%   \keywords, \bmhead)을 쓴다. 블록 1-4 는 원고 조립 때 아래 표시한 자리로 옮긴다.
%   블록 1: \begin{document} 뒤, \maketitle 앞
%   블록 2: Methods 마지막 소절(M15). methods.tex 에 같은 소절이 있으면 하나만 남긴다.
%           M14 Use of large language models 는 methods.tex 의 몫이고 문구는
%           [DECISION: LLM 사용 기재 문구] 다(이 파일에는 본문을 쓰지 않았다).
%   블록 3: 본문 끝, \bibliography 앞(Sci Rep 점검표: Data availability 는 References 앞)
%   블록 4: \bibliography 뒤(\backmatter). Acknowledgements, Author contributions, Additional information
% 인용 키: paper/manuscript/en/refs_front.bib. 키 이름은 다른 절의 FirstAuthorYear 규칙을 따르고,
%   Streletskiy2025 와 munozsabater2021_era5_land 는 refs_intro.bib, refs_results_b.bib 와 같은 키다.
%
% 수치 출처(원고 조립 때 shared/numbers.tex 매크로로 바꾼다. 자릿수는 지침 4.5):
%   -1.6 cm, 95% 신뢰구간 -2.0 to -1.0 .... C1 README 2.3 AB3(-1.63 [-1.98, -1.04], Holm P 0.001)
%   -2.5 cm, -3.0 to -1.9, 1지역 ......... C2 README 2.1 A6(AB4 -2.45 [-3.04, -1.88]), 지역 행 C3 README E36a
%   잔차 ML 추가 차이 확인 못함 ............ C2 README 2.1 A7(AB5 -0.18, Holm P 0.136, 초록 규칙 (b))
%   -0.87, -1.3 cm, 2지역, 캐나다 ......... C5 README 2.1 E2, E3, 2.2 E6, E7(WF4-a, 결과 열람 뒤 설계)
%   13.9, 14.4 cm(라벨 1000개) ............. C4 README 2.1 E1.R1.n1000(13.87 / 14.41)
%   약 99% ................................ C8 README 2.2 B4(격자 안 몫 0.9%, 알래스카 지역 내, 라벨 전량)
%   -2.5, +1.7 cm(라벨 100개) .............. C6 README 2.1 A3, A17(WF2-a 서술 행)
%   다섯 지역, 100 km, 라벨 0-전량 ......... D README 1.3, E24, E25
%   원천 CSV 로 다시 확인한 값(2026-10-04): lgw_bundle.csv(AB3 -1.629003, AB4 -2.454509, AB5 -0.177787),
%   wf_tests.csv(WF4-a -0.873044, -1.320595; WF2-a -2.538178, +1.651555), wf2b_tests.csv(13.86715, 14.40761),
%   wf3b_decomp.csv(share_gain_within 0.008612).
%
% 이 파일의 자리표시: [DECISION: ...] 사용자 결정, [미확인: ...] 확인하지 못한 사실,
%   [XC: ...] [XE: ...] 새 실험(명세 11절 등록부의 문자열 그대로). 제출 전 점검에서 0건이어야 한다.
% 빌드 주의: 자리표시에 한글이 있어 자리표시가 남은 동안은 xelatex(kotex)로 조립한다.
%   pdflatex 빌드는 자리표시를 모두 바꾼 뒤에 한다(지침 M-11).

% 조립(2026-10-05): main.tex 가 \frontblock 을 1, 3, 4 로 정한 뒤 이 파일을 블록마다 \input 한다.
%   각 블록은 \ifnum\frontblock=N ... \fi 로 감쌌다(0 이면 아무것도 내지 않는다). 블록 2 는
%   methods.tex 의 Code availability(M15)와 같은 소절이라 main.tex 가 부르지 않는다(위 블록 2 지시: 하나만 남긴다).
\providecommand{\frontblock}{0}

% =====================================================================
% 블록 1. 제목 쪽, 초록, 핵심어
% =====================================================================
\ifnum\frontblock=1\relax

% 임시 제목(사용자 결정 2026-10-04: 대회 제목의 영역, 나중에 바꾼다). 14단어, 콜론·의문문·금지어 없음.
% [DECISION: 최종 제목]
\title{Permafrost active-layer thickness prediction under sparse observations using physics-based pseudo-label augmentation and machine learning}

% 저자 후보 기록(명세 13.1 의 3번): 제1 저자 후보 백승원(Seungwon Baek), KOPRI 공동 연구 가능자
% 김승희(Seung Hee Kim), 정윤택(Yoon Taek Jung). 저자 자격은 이 파일이 정하지 않는다.
% 형식(템플릿): \author*[1]{\fnm{이름} \sur{성}}\email{...}, \affil*[1]{\orgdiv{}, \orgname{}, \orgaddress{...}}
\author*[1]{\fnm{[DECISION: 저자, 소속, 교신 저자]}}\email{[DECISION: 교신 저자 전자우편]}

\affil*[1]{\orgname{[DECISION: 저자, 소속, 교신 저자]}}

% 초록: 판 B(사용자 잠정 선택, WF 결과는 판정어 없는 수치). 명세 2.2 문안에서 바꾼 곳은 셋이다.
%   (1) 약어 첫 등장 풀이(지침 M-07): RMSE 를 S3 에서 풀어 쓰고, '95% CI' 를 '95% confidence interval' 로 썼다.
%   (2) WF 결과 앞에 '결과 열람 뒤 설계' 표지를 넣었다(명세 1.5, 지침 4.5 끝 항목).
%   (3) 200단어를 지키려고 S2('ML models ... without recalibrated baselines'), S3('Stefan models'),
%       S7(지역을 괄호로), S8('Sub-grid covariates are needed next.')을 줄였다.
% 단어 수(wc -w, 자리표시 제외): 199. 판정어는 S4(AB3 규칙 (a), AB4 규칙 (a)와 지역 수, AB5 규칙 (b))에만 있다.
% [DECISION: 초록 판 A 또는 B]  [DECISION: 판 B 의 WF 수치 개수 허용 여부]
% S8 의 두 자리표시는 결과가 나오면 FINAL_BATCH 의 해당 문장으로 바꾸고, XC 수치를 넣으면 S7 이나 S8 을 대신해 200단어를 지킨다.
\abstract{We provide a label-count workflow for mapping active-layer thickness in sparsely measured permafrost regions with the Stefan model and machine learning (ML). ML models are mostly tested within training regions at one label count, without recalibrated baselines. We held out five regions with 100~km buffers and compared root-mean-square error (RMSE) from zero to all labels against Stefan models with source and recalibrated coefficients (internally pre-registered). Without labels, Stefan pseudo-labels gave lower error than shuffled ones ($-1.6$~cm; 95\% confidence interval $-2.0$ to $-1.0$); with ten labels, recalibration gave lower four-region mean error than the source coefficient ($-2.5$~cm; $-3.0$ to $-1.9$; lower in one region), with no further difference established for residual ML. In analyses designed after these results, method selection by within-target cross-validation with 40 and 160 labels changed RMSE by $-0.87$ and $-1.3$~cm against a fixed residual recipe (two regions, change from Canada). Within Alaska, residual ML reached 13.9 against 14.4~cm RMSE for recalibrated Stefan (1000 labels); about 99\% of the squared-error reduction (all labels) lay between climate grid cells. Spreading 100 labels over blocks changed RMSE against random sampling by $-2.5$~cm (Canada) and $+1.7$~cm (Lena Delta). Sub-grid covariates are needed next. \textcolor{blue}{[PENDING: XE-a, XE-b, satellite inputs (stage 2); registered deadline 8 October 2026, 14:22 KST]}}
% [DECISION: abstract XC sentence] XC is left out of the abstract (FINAL_BATCH 7.1 allows numbers only or exclusion; a numbers-only
% sentence replaces S7 (19 words) and makes the abstract 214 words). Candidate (34 words; would need 14 further words cut):
%   In new random splits of the same regions, the workflow changed RMSE by -1.51 cm against recalibrated Stefan at 160 labels
%   (pooled across regions; over 0.5 cm larger error in 2 of 3 regions).

% 비교판(판 A, WF 제외). 명세 2.3 문안에 같은 약어 수정('95% confidence interval')만 했다. 197단어.
% 판 A 를 고르면 위 \abstract{} 의 내용을 아래 문단으로 바꾼다. 판정어 근거는 명세 2.3 끝 문단.
%   We provide a label-count workflow for mapping active-layer thickness (ALT) in sparsely measured
%   permafrost regions with the Stefan thaw model and machine learning (ML). ML models of ALT are usually
%   evaluated within training regions, at one label count and without a recalibrated physical baseline.
%   We held out four Arctic regions with 100~km buffers and compared errors from zero to all target labels
%   against the Stefan model with source and recalibrated coefficients, in an internally pre-registered
%   analysis. Without labels, Stefan pseudo-labels gave lower error than shuffled ones ($-1.6$~cm; 95\%
%   confidence interval $-2.0$ to $-1.0$), and no difference was established between a physics-model
%   ensemble and the source-coefficient model. With ten labels, coefficient recalibration gave lower
%   four-region mean error ($-2.5$~cm; $-3.0$ to $-1.9$; one of four regions), no further difference was
%   established for residual ML, and residual learning gave lower error than physics outputs used as ML
%   inputs. With all labels, residual ML gave lower four-region mean error than the source-coefficient
%   model, mainly from one region, and lower error than the recalibrated model by less than 0.5~cm. The
%   workflow therefore uses physics at every label count and assigns ML a small correction once labels
%   allow recalibration.

% 핵심어 6개(Sci Rep 상한 6). 명세 1.3 제안. [DECISION: 핵심어]
\keywords{Active-layer thickness, Permafrost, Stefan model, Physics-guided machine learning, Spatial transfer, Sampling design}

\fi

% =====================================================================
% 블록 2. Code availability(Methods 마지막 소절, 명세 6.15 M15)
% =====================================================================
\ifnum\frontblock=2\relax
% 코드 서체(\texttt)는 이 소절의 실제 저장소 경로에만 쓴다(지침 4.1, M-10).
% 경로는 2026-10-04 에 저장소에서 확인했다. 패키지 판은 FINAL_BATCH 1절 '패키지 판' 행,
%   이전 실행의 판은 방법 초안 Table 4(Supplementary Table S16 으로 옮길 예정).
% 저장소 후보: git 원격 git@github.com:01willy/Polar_Bigdata(현재 공개 여부와 공개 시점은 결정 대기).
% TabPFN 판 8.0.7 의 Prior Labs License(v1.2) 10항은 가중치나 원본을 배포할 때 'Built with PriorLabs-TabPFN'
%   표시를 요구한다. 가중치를 재배포하지 않으므로 본문에는 넣지 않았다. 저장소 README 에 넣을지는 공개 때 판단한다.
\subsection*{Code availability}

The code for data assembly, experiments, aggregation, figures and tests is available at [DECISION: 공개 저장소 주소와 공개 시점] and will be archived under the release tag of the reported analyses at [DECISION: Zenodo DOI]. The experiment scripts under \texttt{scripts/} are \texttt{h40\_label\_grid.py} and \texttt{h42\_label\_grid\_ext.py} (transfer), \texttt{h41\_validation\_ladder.py} (validation ladder), \texttt{h54\_workflow.py} (within-region tests), \texttt{h39\_scenarios.py} (aggregation), \texttt{wf0\_reanalysis.py} (post hoc reanalysis) and \texttt{x\_*.py} (additional analyses). The repository also holds the unit tests (\texttt{tests/}), the cloud job configurations (\texttt{configs/rescale/}) and the registration documents with their revision histories. Each output file of a run records the code hash, the configuration hash and the package versions, and the aggregate metadata record the git commit. The cloud runs of the additional analyses use CatBoost 1.2.10, scikit-learn 1.9.1, pandas 2.3.3, SciPy 1.17.1 and NumPy 2.1.1; the versions of earlier runs are listed in Supplementary Table S16. The pretrained weights of TabPFN v2 and TabICL v2 are not redistributed and are obtained from their providers.

\fi

% =====================================================================
% 블록 3. Data availability(본문 끝, References 앞, 명세 6.16)
% =====================================================================
\ifnum\frontblock=3\relax
% 확인 기록(2026-10-04):
%   - 자료 인용 19건은 방법 초안 References 표(67행)와 references/INDEX.md 12.4 표에 있고,
%     DOI 와 제목·연도·저장소를 DataCite API 로 다시 대조했다(refs_front.bib 각 항목 주석).
%   - 새 지역 자료원의 약관은 data/processed/fidelity_base_v4_meta.json 의 by_source 를 따랐다.
%     약관 확인분(license_attention False): calm_pangaea_v4add, firealt, Disko, Ilulissat, myThaw,
%     Palmtag pedon, Du(Zenodo, MIT), Fu(figshare, 지온 유도, 라벨 정의 시험), Kolyma, Indigirka,
%     Mamontov Klyk, Syrdakh, Tavvavuoma, Yamal(PANGAEA 행). 약관 미확인: CALM 누리집 하위 지점(Abisko,
%     Kapp Linné), CUSP v1.1, NSIDC GGD353, Yamal 보고서 유래 행, NSIDC GGD402(기술 통계만).
%     약관 미확인 자료의 인용(Åkerman 1998, Åkerman and Johansson 2008, Jorgenson and Kanevskiy 2025,
%     Petrone 2016, Nixon 2003)은 Supplementary Table S17 과 SI 참고문헌에 둔다. [DECISION: 참고문헌 처리]
%   - CCI ALT v4.0 의 DOI 는 DataCite 에서 제목 'Permafrost active layer thickness for the Northern
%     Hemisphere, v4.0'(2024, CEDA)으로 확인했다. INDEX 01절 esa_cci_permafrost_v4 의 DOI(7479606004...)는
%     DataCite 에 없다. 명세 9.3 의 불일치는 ALT 자료 인용 쪽에서는 d34330ce... 로 정리된다.
%   - SAR: 10.3334/ORNLDAAC/2004 는 DataCite 제목이 'ABoVE: Active Layer Thickness from Airborne L- and
%     P- band SAR, Alaska, 2017, Ver. 3'(Chen et al., 2022)이고, 내려받기 스크립트
%     scripts/0_download/resalt_insar.py 가 이 DOI 로 받은 파일(PDO_ReSALT_*_2017_03.nc4)이 명세의 'ReSALT' 다.
%     PolSAR 상향 ALT(data/raw/polsar_alt/upscaled_alt_{2014,2015,2017}.tif)의 출처는 내려받기 기록이 없다.
%     DataCite 10.3334/ORNLDAAC/2332 의 제목('ABoVE: Upscaled Active Layer Thickness in Northern Alaska,
%     2014-2017', Whitcomb et al., 2025)이 같은 자료로 보이나 확인하지 못했다.
%   - Copernicus DEM DOI 10.5270/ESA-c5d3d65 는 doi.org 에서 COP-DEM 자료 설명 쪽으로 연결된다(인용 형식은 미확인).
%   - 지도 마스크(CCI 영구동토 범위 평균, JRC 지표수)는 방법 초안 Data availability 의 문장을 따랐다.
\section*{Data availability}

The active-layer labels of the main regions are available from the GTN-P/CALM compilation\cite{Streletskiy2025} (PANGAEA, \url{https://doi.org/10.1594/PANGAEA.972777}), the ALLena compilation for the Lena River Delta\cite{Veremeeva2025} (PANGAEA, \url{https://doi.org/10.1594/PANGAEA.973813}) and the ABoVE soil moisture and active-layer thickness data set, version~2\cite{Moore2025} (ORNL DAAC, \url{https://doi.org/10.3334/ORNLDAAC/2369}). The labels of the additional regions come from published data sets\cite{DuE2026,Grosse2007,Palmtag2022,Talucci2024,Pohl2026,Makarieva2017,Liang2023,Walker2009,Scheer2024,Martin2023,Boike2024,Hammar2025,Zastruzny2024,Sannel2020}, and the temperature-derived labels used only in the label-definition test come from a Tibetan Plateau compilation\cite{Fu2025}. Their identifiers and licences are listed in Supplementary Table S17. Labels from sources without a confirmed redistribution licence were excluded from the licence-verified edition used for the main results and are not redistributed; they can be obtained from the original providers listed in Supplementary Table S17. The full North Atlantic edition, which includes such labels, is reported in the Supplementary Information only.

The covariates were derived from ERA5-Land monthly averaged data\cite{munozsabater2021_era5_land} (Copernicus Climate Data Store, \url{https://doi.org/10.24381/cds.68d2bb30}), SoilGrids 2.0\cite{Poggio2021} [미확인: 접근 시점의 SoilGrids 판 표지], the Copernicus DEM GLO-30 (\url{https://doi.org/10.5270/ESA-c5d3d65}) [미확인: Copernicus DEM 인용 형식] and the ESA CCI Permafrost active-layer thickness product, version 4.0\cite{Westermann2024} (CEDA, \url{https://doi.org/10.5285/d34330ce3f604e368c06d76de1987ce5}). The synthetic-aperture radar (SAR) covariates of the within-Alaska test came from the ABoVE airborne L- and P-band SAR active-layer retrievals for 2017 (ORNL DAAC, \url{https://doi.org/10.3334/ORNLDAAC/2004}) and from upscaled P-band polarimetric SAR active-layer maps of northern Alaska\cite{Whitcomb2024} [미확인: PolSAR 상향 ALT 자료의 DOI, 후보 10.3334/ORNLDAAC/2332]. The map masks use the ESA CCI Permafrost extent product [미확인: 영구동토 범위 자료의 DOI와 판] and the JRC Global Surface Water occurrence layer\cite{Pekel2016}.

The assembled label and covariate tables, the run outputs, the aggregate and verdict tables, the cell-level predictions and the source data of all figures will be deposited at [DECISION: 저장소와 DOI]. Rows from sources without a confirmed licence are excluded from the deposited tables. [DECISION: KPDC 자료, SI 포함 여부와 식별자, 이용 조건]

\fi

% =====================================================================
% 블록 4. \bibliography 뒤(\backmatter)
% =====================================================================
\ifnum\frontblock=4\relax
% Sci Rep: 감사의 글은 짧게, 저자 기여는 참고문헌 뒤에 모든 저자별로, 이해상충은 'Additional information' 아래.
% 세 문구 모두 [DECISION: Acknowledgements, Author contributions, Competing interests 문구](명세 13.1 의 5번).

\bmhead{Acknowledgements}

[DECISION: Acknowledgements, Author contributions, Competing interests 문구] We thank the providers of the data sets listed under Data availability. The active-layer thickness covariate uses data of the ESA Climate Change Initiative Permafrost\_cci project. Terrain covariates were produced using Copernicus WorldDEM-30 \textcopyright{} DLR e.V. 2010--2014 and \textcopyright{} Airbus Defence and Space GmbH 2014--2018, provided under COPERNICUS by the European Union and ESA; all rights reserved. [미확인: Copernicus DEM 사용 허가가 요구하는 귀속 문구의 원문 대조] [미확인: ERA5-Land(Copernicus Climate Change Service) 사용 허가가 요구하는 귀속 문구] [DECISION: KPDC 자료, SI 포함 여부와 식별자, 이용 조건] [DECISION: 대회 보고서 인용 여부] [DECISION: 연구비 지원 기관과 과제 번호, 또는 지원 없음]

% KPDC 자료를 SI 에 두면 KPDC 감사 문구를 더한다(앞부분 초안 3절의 문안. 식별자 2707, 2955 확인 뒤).
% 대회: '2026 극지 빅데이터-인공지능 활용 경진대회'(주최 극지연구소, 주관 KPDC). 공식 영문 명칭 [미확인].
%   선행 공개 사실은 표지 편지에 적고 사본을 첨부한다(명세 13.1 의 12번).

\bmhead{Author contributions}

[DECISION: Acknowledgements, Author contributions, Competing interests 문구] [DECISION: 저자, 소속, 교신 저자]

% 형식(결정 뒤): 모든 저자에 대해 CRediT 역할(Conceptualization, Methodology, Software, Validation,
%   Formal analysis, Investigation, Resources, Data curation, Writing (original draft), Writing (review and
%   editing), Visualization, Supervision, Project administration, Funding acquisition)을 이름과 함께 적고,
%   끝에 'All authors reviewed the manuscript.' 같은 문장을 둔다. 저자 자격은 이 파일이 정하지 않는다.

\bmhead{Additional information}

Competing interests: [DECISION: Acknowledgements, Author contributions, Competing interests 문구]

Correspondence and requests for materials should be addressed to [DECISION: 저자, 소속, 교신 저자].

% Competing interests 문안(결정 뒤, 해당 없을 때 Sci Rep 예시 문구): 'The authors declare no competing interests.'
\fi


\bibliography{references}

\backmatter

% Block 4 of front.tex: Acknowledgements, Author contributions, Additional information
\def\frontblock{4}% front.tex : 영문 원고의 제목 쪽, 초록, 핵심어, Code availability, Data availability,
%             Acknowledgements, Author contributions, Additional information(Competing interests)
% 작성 2026-10-04. LaTeX 본문 조각이다(프리앰블 없음). 명세: paper/manuscript/MANUSCRIPT_SPEC.md
%   1.1(순서), 1.3(제목·핵심어), 1.5(서술 방향), 2.1-2.4(초록), 6.15(M15 Code availability),
%   6.16(Data availability), 13.1(사용자 결정 대기).
% 서식: Springer Nature 템플릿 sn-jnl(sn-nature 선택지)의 명령(\title, \author, \affil, \abstract,
%   \keywords, \bmhead)을 쓴다. 블록 1-4 는 원고 조립 때 아래 표시한 자리로 옮긴다.
%   블록 1: \begin{document} 뒤, \maketitle 앞
%   블록 2: Methods 마지막 소절(M15). methods.tex 에 같은 소절이 있으면 하나만 남긴다.
%           M14 Use of large language models 는 methods.tex 의 몫이고 문구는
%           [DECISION: LLM 사용 기재 문구] 다(이 파일에는 본문을 쓰지 않았다).
%   블록 3: 본문 끝, \bibliography 앞(Sci Rep 점검표: Data availability 는 References 앞)
%   블록 4: \bibliography 뒤(\backmatter). Acknowledgements, Author contributions, Additional information
% 인용 키: paper/manuscript/en/refs_front.bib. 키 이름은 다른 절의 FirstAuthorYear 규칙을 따르고,
%   Streletskiy2025 와 munozsabater2021_era5_land 는 refs_intro.bib, refs_results_b.bib 와 같은 키다.
%
% 수치 출처(원고 조립 때 shared/numbers.tex 매크로로 바꾼다. 자릿수는 지침 4.5):
%   -1.6 cm, 95% 신뢰구간 -2.0 to -1.0 .... C1 README 2.3 AB3(-1.63 [-1.98, -1.04], Holm P 0.001)
%   -2.5 cm, -3.0 to -1.9, 1지역 ......... C2 README 2.1 A6(AB4 -2.45 [-3.04, -1.88]), 지역 행 C3 README E36a
%   잔차 ML 추가 차이 확인 못함 ............ C2 README 2.1 A7(AB5 -0.18, Holm P 0.136, 초록 규칙 (b))
%   -0.87, -1.3 cm, 2지역, 캐나다 ......... C5 README 2.1 E2, E3, 2.2 E6, E7(WF4-a, 결과 열람 뒤 설계)
%   13.9, 14.4 cm(라벨 1000개) ............. C4 README 2.1 E1.R1.n1000(13.87 / 14.41)
%   약 99% ................................ C8 README 2.2 B4(격자 안 몫 0.9%, 알래스카 지역 내, 라벨 전량)
%   -2.5, +1.7 cm(라벨 100개) .............. C6 README 2.1 A3, A17(WF2-a 서술 행)
%   다섯 지역, 100 km, 라벨 0-전량 ......... D README 1.3, E24, E25
%   원천 CSV 로 다시 확인한 값(2026-10-04): lgw_bundle.csv(AB3 -1.629003, AB4 -2.454509, AB5 -0.177787),
%   wf_tests.csv(WF4-a -0.873044, -1.320595; WF2-a -2.538178, +1.651555), wf2b_tests.csv(13.86715, 14.40761),
%   wf3b_decomp.csv(share_gain_within 0.008612).
%
% 이 파일의 자리표시: [DECISION: ...] 사용자 결정, [미확인: ...] 확인하지 못한 사실,
%   [XC: ...] [XE: ...] 새 실험(명세 11절 등록부의 문자열 그대로). 제출 전 점검에서 0건이어야 한다.
% 빌드 주의: 자리표시에 한글이 있어 자리표시가 남은 동안은 xelatex(kotex)로 조립한다.
%   pdflatex 빌드는 자리표시를 모두 바꾼 뒤에 한다(지침 M-11).

% 조립(2026-10-05): main.tex 가 \frontblock 을 1, 3, 4 로 정한 뒤 이 파일을 블록마다 \input 한다.
%   각 블록은 \ifnum\frontblock=N ... \fi 로 감쌌다(0 이면 아무것도 내지 않는다). 블록 2 는
%   methods.tex 의 Code availability(M15)와 같은 소절이라 main.tex 가 부르지 않는다(위 블록 2 지시: 하나만 남긴다).
\providecommand{\frontblock}{0}

% =====================================================================
% 블록 1. 제목 쪽, 초록, 핵심어
% =====================================================================
\ifnum\frontblock=1\relax

% 임시 제목(사용자 결정 2026-10-04: 대회 제목의 영역, 나중에 바꾼다). 14단어, 콜론·의문문·금지어 없음.
% [DECISION: 최종 제목]
\title{Permafrost active-layer thickness prediction under sparse observations using physics-based pseudo-label augmentation and machine learning}

% 저자 후보 기록(명세 13.1 의 3번): 제1 저자 후보 백승원(Seungwon Baek), KOPRI 공동 연구 가능자
% 김승희(Seung Hee Kim), 정윤택(Yoon Taek Jung). 저자 자격은 이 파일이 정하지 않는다.
% 형식(템플릿): \author*[1]{\fnm{이름} \sur{성}}\email{...}, \affil*[1]{\orgdiv{}, \orgname{}, \orgaddress{...}}
\author*[1]{\fnm{[DECISION: 저자, 소속, 교신 저자]}}\email{[DECISION: 교신 저자 전자우편]}

\affil*[1]{\orgname{[DECISION: 저자, 소속, 교신 저자]}}

% 초록: 판 B(사용자 잠정 선택, WF 결과는 판정어 없는 수치). 명세 2.2 문안에서 바꾼 곳은 셋이다.
%   (1) 약어 첫 등장 풀이(지침 M-07): RMSE 를 S3 에서 풀어 쓰고, '95% CI' 를 '95% confidence interval' 로 썼다.
%   (2) WF 결과 앞에 '결과 열람 뒤 설계' 표지를 넣었다(명세 1.5, 지침 4.5 끝 항목).
%   (3) 200단어를 지키려고 S2('ML models ... without recalibrated baselines'), S3('Stefan models'),
%       S7(지역을 괄호로), S8('Sub-grid covariates are needed next.')을 줄였다.
% 단어 수(wc -w, 자리표시 제외): 199. 판정어는 S4(AB3 규칙 (a), AB4 규칙 (a)와 지역 수, AB5 규칙 (b))에만 있다.
% [DECISION: 초록 판 A 또는 B]  [DECISION: 판 B 의 WF 수치 개수 허용 여부]
% S8 의 두 자리표시는 결과가 나오면 FINAL_BATCH 의 해당 문장으로 바꾸고, XC 수치를 넣으면 S7 이나 S8 을 대신해 200단어를 지킨다.
\abstract{We provide a label-count workflow for mapping active-layer thickness in sparsely measured permafrost regions with the Stefan model and machine learning (ML). ML models are mostly tested within training regions at one label count, without recalibrated baselines. We held out five regions with 100~km buffers and compared root-mean-square error (RMSE) from zero to all labels against Stefan models with source and recalibrated coefficients (internally pre-registered). Without labels, Stefan pseudo-labels gave lower error than shuffled ones ($-1.6$~cm; 95\% confidence interval $-2.0$ to $-1.0$); with ten labels, recalibration gave lower four-region mean error than the source coefficient ($-2.5$~cm; $-3.0$ to $-1.9$; lower in one region), with no further difference established for residual ML. In analyses designed after these results, method selection by within-target cross-validation with 40 and 160 labels changed RMSE by $-0.87$ and $-1.3$~cm against a fixed residual recipe (two regions, change from Canada). Within Alaska, residual ML reached 13.9 against 14.4~cm RMSE for recalibrated Stefan (1000 labels); about 99\% of the squared-error reduction (all labels) lay between climate grid cells. Spreading 100 labels over blocks changed RMSE against random sampling by $-2.5$~cm (Canada) and $+1.7$~cm (Lena Delta). Sub-grid covariates are needed next. \textcolor{blue}{[PENDING: XE-a, XE-b, satellite inputs (stage 2); registered deadline 8 October 2026, 14:22 KST]}}
% [DECISION: abstract XC sentence] XC is left out of the abstract (FINAL_BATCH 7.1 allows numbers only or exclusion; a numbers-only
% sentence replaces S7 (19 words) and makes the abstract 214 words). Candidate (34 words; would need 14 further words cut):
%   In new random splits of the same regions, the workflow changed RMSE by -1.51 cm against recalibrated Stefan at 160 labels
%   (pooled across regions; over 0.5 cm larger error in 2 of 3 regions).

% 비교판(판 A, WF 제외). 명세 2.3 문안에 같은 약어 수정('95% confidence interval')만 했다. 197단어.
% 판 A 를 고르면 위 \abstract{} 의 내용을 아래 문단으로 바꾼다. 판정어 근거는 명세 2.3 끝 문단.
%   We provide a label-count workflow for mapping active-layer thickness (ALT) in sparsely measured
%   permafrost regions with the Stefan thaw model and machine learning (ML). ML models of ALT are usually
%   evaluated within training regions, at one label count and without a recalibrated physical baseline.
%   We held out four Arctic regions with 100~km buffers and compared errors from zero to all target labels
%   against the Stefan model with source and recalibrated coefficients, in an internally pre-registered
%   analysis. Without labels, Stefan pseudo-labels gave lower error than shuffled ones ($-1.6$~cm; 95\%
%   confidence interval $-2.0$ to $-1.0$), and no difference was established between a physics-model
%   ensemble and the source-coefficient model. With ten labels, coefficient recalibration gave lower
%   four-region mean error ($-2.5$~cm; $-3.0$ to $-1.9$; one of four regions), no further difference was
%   established for residual ML, and residual learning gave lower error than physics outputs used as ML
%   inputs. With all labels, residual ML gave lower four-region mean error than the source-coefficient
%   model, mainly from one region, and lower error than the recalibrated model by less than 0.5~cm. The
%   workflow therefore uses physics at every label count and assigns ML a small correction once labels
%   allow recalibration.

% 핵심어 6개(Sci Rep 상한 6). 명세 1.3 제안. [DECISION: 핵심어]
\keywords{Active-layer thickness, Permafrost, Stefan model, Physics-guided machine learning, Spatial transfer, Sampling design}

\fi

% =====================================================================
% 블록 2. Code availability(Methods 마지막 소절, 명세 6.15 M15)
% =====================================================================
\ifnum\frontblock=2\relax
% 코드 서체(\texttt)는 이 소절의 실제 저장소 경로에만 쓴다(지침 4.1, M-10).
% 경로는 2026-10-04 에 저장소에서 확인했다. 패키지 판은 FINAL_BATCH 1절 '패키지 판' 행,
%   이전 실행의 판은 방법 초안 Table 4(Supplementary Table S16 으로 옮길 예정).
% 저장소 후보: git 원격 git@github.com:01willy/Polar_Bigdata(현재 공개 여부와 공개 시점은 결정 대기).
% TabPFN 판 8.0.7 의 Prior Labs License(v1.2) 10항은 가중치나 원본을 배포할 때 'Built with PriorLabs-TabPFN'
%   표시를 요구한다. 가중치를 재배포하지 않으므로 본문에는 넣지 않았다. 저장소 README 에 넣을지는 공개 때 판단한다.
\subsection*{Code availability}

The code for data assembly, experiments, aggregation, figures and tests is available at [DECISION: 공개 저장소 주소와 공개 시점] and will be archived under the release tag of the reported analyses at [DECISION: Zenodo DOI]. The experiment scripts under \texttt{scripts/} are \texttt{h40\_label\_grid.py} and \texttt{h42\_label\_grid\_ext.py} (transfer), \texttt{h41\_validation\_ladder.py} (validation ladder), \texttt{h54\_workflow.py} (within-region tests), \texttt{h39\_scenarios.py} (aggregation), \texttt{wf0\_reanalysis.py} (post hoc reanalysis) and \texttt{x\_*.py} (additional analyses). The repository also holds the unit tests (\texttt{tests/}), the cloud job configurations (\texttt{configs/rescale/}) and the registration documents with their revision histories. Each output file of a run records the code hash, the configuration hash and the package versions, and the aggregate metadata record the git commit. The cloud runs of the additional analyses use CatBoost 1.2.10, scikit-learn 1.9.1, pandas 2.3.3, SciPy 1.17.1 and NumPy 2.1.1; the versions of earlier runs are listed in Supplementary Table S16. The pretrained weights of TabPFN v2 and TabICL v2 are not redistributed and are obtained from their providers.

\fi

% =====================================================================
% 블록 3. Data availability(본문 끝, References 앞, 명세 6.16)
% =====================================================================
\ifnum\frontblock=3\relax
% 확인 기록(2026-10-04):
%   - 자료 인용 19건은 방법 초안 References 표(67행)와 references/INDEX.md 12.4 표에 있고,
%     DOI 와 제목·연도·저장소를 DataCite API 로 다시 대조했다(refs_front.bib 각 항목 주석).
%   - 새 지역 자료원의 약관은 data/processed/fidelity_base_v4_meta.json 의 by_source 를 따랐다.
%     약관 확인분(license_attention False): calm_pangaea_v4add, firealt, Disko, Ilulissat, myThaw,
%     Palmtag pedon, Du(Zenodo, MIT), Fu(figshare, 지온 유도, 라벨 정의 시험), Kolyma, Indigirka,
%     Mamontov Klyk, Syrdakh, Tavvavuoma, Yamal(PANGAEA 행). 약관 미확인: CALM 누리집 하위 지점(Abisko,
%     Kapp Linné), CUSP v1.1, NSIDC GGD353, Yamal 보고서 유래 행, NSIDC GGD402(기술 통계만).
%     약관 미확인 자료의 인용(Åkerman 1998, Åkerman and Johansson 2008, Jorgenson and Kanevskiy 2025,
%     Petrone 2016, Nixon 2003)은 Supplementary Table S17 과 SI 참고문헌에 둔다. [DECISION: 참고문헌 처리]
%   - CCI ALT v4.0 의 DOI 는 DataCite 에서 제목 'Permafrost active layer thickness for the Northern
%     Hemisphere, v4.0'(2024, CEDA)으로 확인했다. INDEX 01절 esa_cci_permafrost_v4 의 DOI(7479606004...)는
%     DataCite 에 없다. 명세 9.3 의 불일치는 ALT 자료 인용 쪽에서는 d34330ce... 로 정리된다.
%   - SAR: 10.3334/ORNLDAAC/2004 는 DataCite 제목이 'ABoVE: Active Layer Thickness from Airborne L- and
%     P- band SAR, Alaska, 2017, Ver. 3'(Chen et al., 2022)이고, 내려받기 스크립트
%     scripts/0_download/resalt_insar.py 가 이 DOI 로 받은 파일(PDO_ReSALT_*_2017_03.nc4)이 명세의 'ReSALT' 다.
%     PolSAR 상향 ALT(data/raw/polsar_alt/upscaled_alt_{2014,2015,2017}.tif)의 출처는 내려받기 기록이 없다.
%     DataCite 10.3334/ORNLDAAC/2332 의 제목('ABoVE: Upscaled Active Layer Thickness in Northern Alaska,
%     2014-2017', Whitcomb et al., 2025)이 같은 자료로 보이나 확인하지 못했다.
%   - Copernicus DEM DOI 10.5270/ESA-c5d3d65 는 doi.org 에서 COP-DEM 자료 설명 쪽으로 연결된다(인용 형식은 미확인).
%   - 지도 마스크(CCI 영구동토 범위 평균, JRC 지표수)는 방법 초안 Data availability 의 문장을 따랐다.
\section*{Data availability}

The active-layer labels of the main regions are available from the GTN-P/CALM compilation\cite{Streletskiy2025} (PANGAEA, \url{https://doi.org/10.1594/PANGAEA.972777}), the ALLena compilation for the Lena River Delta\cite{Veremeeva2025} (PANGAEA, \url{https://doi.org/10.1594/PANGAEA.973813}) and the ABoVE soil moisture and active-layer thickness data set, version~2\cite{Moore2025} (ORNL DAAC, \url{https://doi.org/10.3334/ORNLDAAC/2369}). The labels of the additional regions come from published data sets\cite{DuE2026,Grosse2007,Palmtag2022,Talucci2024,Pohl2026,Makarieva2017,Liang2023,Walker2009,Scheer2024,Martin2023,Boike2024,Hammar2025,Zastruzny2024,Sannel2020}, and the temperature-derived labels used only in the label-definition test come from a Tibetan Plateau compilation\cite{Fu2025}. Their identifiers and licences are listed in Supplementary Table S17. Labels from sources without a confirmed redistribution licence were excluded from the licence-verified edition used for the main results and are not redistributed; they can be obtained from the original providers listed in Supplementary Table S17. The full North Atlantic edition, which includes such labels, is reported in the Supplementary Information only.

The covariates were derived from ERA5-Land monthly averaged data\cite{munozsabater2021_era5_land} (Copernicus Climate Data Store, \url{https://doi.org/10.24381/cds.68d2bb30}), SoilGrids 2.0\cite{Poggio2021} [미확인: 접근 시점의 SoilGrids 판 표지], the Copernicus DEM GLO-30 (\url{https://doi.org/10.5270/ESA-c5d3d65}) [미확인: Copernicus DEM 인용 형식] and the ESA CCI Permafrost active-layer thickness product, version 4.0\cite{Westermann2024} (CEDA, \url{https://doi.org/10.5285/d34330ce3f604e368c06d76de1987ce5}). The synthetic-aperture radar (SAR) covariates of the within-Alaska test came from the ABoVE airborne L- and P-band SAR active-layer retrievals for 2017 (ORNL DAAC, \url{https://doi.org/10.3334/ORNLDAAC/2004}) and from upscaled P-band polarimetric SAR active-layer maps of northern Alaska\cite{Whitcomb2024} [미확인: PolSAR 상향 ALT 자료의 DOI, 후보 10.3334/ORNLDAAC/2332]. The map masks use the ESA CCI Permafrost extent product [미확인: 영구동토 범위 자료의 DOI와 판] and the JRC Global Surface Water occurrence layer\cite{Pekel2016}.

The assembled label and covariate tables, the run outputs, the aggregate and verdict tables, the cell-level predictions and the source data of all figures will be deposited at [DECISION: 저장소와 DOI]. Rows from sources without a confirmed licence are excluded from the deposited tables. [DECISION: KPDC 자료, SI 포함 여부와 식별자, 이용 조건]

\fi

% =====================================================================
% 블록 4. \bibliography 뒤(\backmatter)
% =====================================================================
\ifnum\frontblock=4\relax
% Sci Rep: 감사의 글은 짧게, 저자 기여는 참고문헌 뒤에 모든 저자별로, 이해상충은 'Additional information' 아래.
% 세 문구 모두 [DECISION: Acknowledgements, Author contributions, Competing interests 문구](명세 13.1 의 5번).

\bmhead{Acknowledgements}

[DECISION: Acknowledgements, Author contributions, Competing interests 문구] We thank the providers of the data sets listed under Data availability. The active-layer thickness covariate uses data of the ESA Climate Change Initiative Permafrost\_cci project. Terrain covariates were produced using Copernicus WorldDEM-30 \textcopyright{} DLR e.V. 2010--2014 and \textcopyright{} Airbus Defence and Space GmbH 2014--2018, provided under COPERNICUS by the European Union and ESA; all rights reserved. [미확인: Copernicus DEM 사용 허가가 요구하는 귀속 문구의 원문 대조] [미확인: ERA5-Land(Copernicus Climate Change Service) 사용 허가가 요구하는 귀속 문구] [DECISION: KPDC 자료, SI 포함 여부와 식별자, 이용 조건] [DECISION: 대회 보고서 인용 여부] [DECISION: 연구비 지원 기관과 과제 번호, 또는 지원 없음]

% KPDC 자료를 SI 에 두면 KPDC 감사 문구를 더한다(앞부분 초안 3절의 문안. 식별자 2707, 2955 확인 뒤).
% 대회: '2026 극지 빅데이터-인공지능 활용 경진대회'(주최 극지연구소, 주관 KPDC). 공식 영문 명칭 [미확인].
%   선행 공개 사실은 표지 편지에 적고 사본을 첨부한다(명세 13.1 의 12번).

\bmhead{Author contributions}

[DECISION: Acknowledgements, Author contributions, Competing interests 문구] [DECISION: 저자, 소속, 교신 저자]

% 형식(결정 뒤): 모든 저자에 대해 CRediT 역할(Conceptualization, Methodology, Software, Validation,
%   Formal analysis, Investigation, Resources, Data curation, Writing (original draft), Writing (review and
%   editing), Visualization, Supervision, Project administration, Funding acquisition)을 이름과 함께 적고,
%   끝에 'All authors reviewed the manuscript.' 같은 문장을 둔다. 저자 자격은 이 파일이 정하지 않는다.

\bmhead{Additional information}

Competing interests: [DECISION: Acknowledgements, Author contributions, Competing interests 문구]

Correspondence and requests for materials should be addressed to [DECISION: 저자, 소속, 교신 저자].

% Competing interests 문안(결정 뒤, 해당 없을 때 Sci Rep 예시 문구): 'The authors declare no competing interests.'
\fi


\clearpage
% legends.tex : Figure legends (Fig. 1 to 7) and the Table 1 caption.
% Body text only, no preamble. Sci Rep places legends after the references.
% Spec: paper/manuscript/MANUSCRIPT_SPEC.md sections 7 and 11; panel content and
% first sentences from the v3 figure spec (v3_restructure/FIGURE_SPEC_v3.md)
% (legend drafts in sections 3.6 to 10.4). First sentences are kept as in the spec.
% Limits: each figure legend <= 350 words after the X placeholder is replaced
% (spec budget: Fig 1, 4, 6, 7 reserve 40 to 50 words for the placeholder);
% Table 1 caption <= 150 words. Count on the rendered PDF (pdftotext, wc -w).
% Numbers: read from the source tables named in paper/claims/*/README.md and
% rounded with style guide 4.5 from the unrounded file values (|delta| < 1 cm two
% decimals, other cm values one decimal, correlations two decimals, percentages
% integer, CI ends with the digits of the point estimate). Replace with
% numbers.tex macros once the number-building script exists (QA item M-09).
% Sign and symbols: U+2212 for negative numbers and U+00B0 for degrees, as in the
% sibling section files; numeric ranges use "--"; ranges with a negative end use
% "to". Panel references use (a), (b); letter ranges use "to".
% Placeholders [X?: ...] are the registered strings of spec section 11, copied
% character for character. They contain Korean text, so a build that keeps them
% needs xeCJK. [미확인: ...] marks facts not yet confirmed (spec, figure spec 14).
% Citation keys: paper/manuscript/en/refs_legends.bib (keys follow references/INDEX.md).

\section*{Figure legends}

% ---- Figure 1 -------------------------------------------------------------
% Sources: D README 2 (licence counts, source coefficient, ERA5-Land 2015-2020
% period, caution 5); figure spec 3.3 (panels, 400-cell sample, white circles, map
% elements); final batch plan 2.8 (validation ladder: three regions, Stefan fitted
% on the training fold). Central meridian of panel a: fig1.py uses the Lena Delta zoom centre
% (LON0_A = None -> lon_c 126.698 E, Fig1_values.txt), not 45 W (checked 2026-10-04).
\noindent Figure 1. Study regions, Stefan coefficients and evaluation design.
(a) Labelled cells aggregated to 0.5° blocks; circle area is proportional to the number of 1~km label locations per block, and white circles mark regions added after the label tables were fixed (Central Russia, expanded edition; Tibetan Plateau, inset; Table~1).
Shading marks continuous (permafrost fraction $\geq$90\%) and discontinuous (50--90\%) permafrost [미확인: 영구동토 구역 자료 이름과 판]; the rectangle marks panels c and d.
(b) Stefan coefficient $E = \mathrm{ALT}/\sqrt{\mathrm{TDD}}$ of each labelled cell on a logarithmic axis, where ALT is the active-layer thickness and TDD the thawing degree-days of 2~m air temperature (ERA5-Land\cite{munozsabater2021_era5_land}, 2015--2020 mean); at most 400 cells are shown per region, and Source coefficient marks the coefficient $E_0$ fitted on the source pool when the region is held out.
(c) Region holdout: the Lena Delta cells and all cells within 100~km of them are removed from the training data.
(d) Label blocks and scoring blocks of one of five splits, with 40 label cells drawn from the label blocks; errors are scored only in scoring blocks.
Within-region tests use such splits without source data and were designed after the transfer results had been viewed.
(e) Root-mean-square error (RMSE) of direct machine learning (ML) and of the Stefan model fitted to the training cells (Recalibrated Stefan) under random-cell, site, 0.5° block, $k$-fold nearest-neighbour distance matching (kNNDM)\cite{linnenbrink2024_knndm} and region-holdout validation, against the median distance from scoring cells to the nearest training cell (mean of Alaska, the Lena Delta and Canada, with regional values); [XH: Fig 1e | 결과 열람 뒤 설계, 서술(판정어 없음) | 2.8 사전 고정 서술 규칙].
Cells without a verified licence (Canada 38, North Atlantic 19, W Russia 1) are not drawn.
Polar stereographic projection (central meridian 127°~E), scale bar true at 70°~N; Tibetan Plateau inset in a Lambert azimuthal equal-area projection; coastlines, Natural Earth 50~m; Cartopy [미확인: 판].

\medskip
% ---- Figure 2 -------------------------------------------------------------
% Sources: C2 README 2.1 (four-region recalibration and residual rows at ten
% labels; uncorrected P 0.034 from the abstract bundle, column p_two), 2.2 (Canada and W Russia rows); C3 README 2.5; figure table
% fig2_region_curves (W Russia residual curve at ten labels, -12.52); resamples
% from the pool curve meta (10,000) and the regional curve table (1000). Both CIs
% agree for the four-region recalibration row and the Canada and W Russia rows.
\noindent Figure 2. With ten labels, Stefan recalibration lowered the four-region mean error; no further difference was established for residual ML.
Error change is the root-mean-square error (RMSE) of a method minus that of the source-coefficient Stefan model (zero line) on the scoring blocks of each held-out region; negative values mean lower error.
Recalibrated Stefan refits the coefficient $E$ with the $n$ target labels, shrunk towards the source coefficient ($\kappa = 10$); Anchor + residual ML adds a CatBoost residual with weight 0.25; Direct ML is CatBoost without a physics anchor; Year-matched Stefan uses thawing degree-days of the label years.
(a) Mean of the four main regions (the Lena Delta, Canada, W Russia, E Russia), up to ten labels.
(b) Mean of the Lena Delta and Canada, the regions with at least 40 candidate labels; All, every label; Year-matched Stefan is a point estimate.
(c to f) Single regions; label counts above the candidate pool (14--17 cells in W and E Russia, 325--406 in Canada) are masked, and panel e has its own vertical range (values to −12.5~cm).
Lines are cell-weighted means over five splits; bands are 95\% block-bootstrap confidence intervals (CIs) of the two recalibration-based methods (10,000 resamples for the fixed-composition means in panels a and b; 1000 in panels c to f).
With ten labels, recalibration lowered the four-region RMSE by 2.5~cm (95\% CI 1.9 to 3.0~cm; Holm-adjusted $P = 0.001$), with lower error in W Russia only (12.0~cm; 95\% CI 10.5 to 13.6~cm) and higher error in Canada (2.3~cm; 95\% CI 0.7 to 2.9~cm); adding the residual changed it by −0.18~cm (95\% CI −0.53 to −0.01~cm; $P = 0.034$ before correction, Holm-adjusted $P = 0.136$).
Verdicts of all plotted contrasts are in Supplementary Table~S2.
The label-grid analyses were internally pre-registered and run on one computing platform.

\medskip
% ---- Figure 3 -------------------------------------------------------------
% Sources: C1 README 2.3 (shuffled and linear placebo rows; abstract-bundle Holm P);
% C3 README 2.1 and 2.2 (structure rows at 0, 3, 10, 40, 160 labels; the 40-label
% physics-output row depends on the split-independence assumption), 2.6 (least
% squares minus shrinkage, descriptive). Plotted series: figure table
% fig3_b_L10_alln (physics outputs, recalibrated inputs). Resamples 10,000.
% Panel a: the rendered figure (fig3.py CONCEPT_REGION, Fig3_values.txt [6]-[7]) uses W Russia (31 cells, E0 1.589510,
% ten-label recalibrated E 2.109580), not the Lena Delta of the spec draft; the legend follows the rendered figure.
\noindent Figure 3. Physics pseudo-labels lowered error relative to shuffled pseudo-labels, and the better combination structure depended on the number of labels.
(a) Concept on the 31 labelled cells of W Russia: active-layer thickness against the square root of thawing degree-days, with the source-coefficient Stefan line (Source Stefan, 1.59~cm~(°C~d)$^{-1/2}$), the line recalibrated with ten drawn labels (Recalibrated Stefan, 2.11) and the departures of those labels that residual ML learns (Residual); thickness increases downwards.
(b) Direct ML trained with physics pseudo-labels minus the same model trained with placebo pseudo-labels, four-region mean, without labels and with ten labels.
(c) Residual structure minus physics-input structure against the number of labels, for two input sets (Physics inputs, outputs of two physics models; Recalibrated inputs, source and recalibrated Stefan predictions); the key separates registered label counts (0 and 10) from auxiliary ones, and values at 40 and 160 labels are Lena Delta and Canada means.
(d) Local least-squares minus shrinkage recalibration of the Stefan coefficient (descriptive; positive values favour shrinkage), four-region mean up to ten labels, Lena Delta and Canada mean beyond.
Intervals are 95\% block-bootstrap confidence intervals (CIs; 10,000 resamples of 0.5° scoring blocks), weighted by cell and by block as indicated in the key (panel d, by cell); the grey band marks ±0.5~cm.
In panel b, physics pseudo-labels lowered error relative to shuffled pseudo-labels by 1.6~cm without labels (95\% CI 1.0 to 2.0~cm; Holm-adjusted $P = 0.001$); the difference from linear thaw-index pseudo-labels was 0.48~cm without labels and within ±0.5~cm with ten labels.
In panel c, the residual structure had 2.5--3.7~cm lower error under both weightings with up to ten labels (Holm-adjusted $P = 0.001$ for physics inputs at ten labels); at 40 and 160 labels the physics-input structure had 1.9--2.4~cm lower error, a difference that came from Canada (one of four contrasts depends on the split-independence assumption).
The placebo and structure analyses were internally pre-registered.

\medskip
% ---- Figure 4 -------------------------------------------------------------
% Sources: C2 README 2.1 (rank correlation 0.38 [0.17, 0.55], 28 targets, signed
% -0.01), 2.2 Tibetan rows (bias 227.71, gain 159.55); C5 README 2.1 and 2.2
% (selection minus fixed recipe at 10, 40, 160 and all labels; regional rows);
% 1.2 N4 (40-label reduction and all-label non-inferiority depend on the
% split-independence assumption). Placeholder: spec section 11.
\noindent Figure 4. The error reduction of label-based correction was rank-correlated with the source-coefficient bias measured with ten labels.
Intervals in panels b and c are 95\% block-bootstrap confidence intervals (CIs; 10,000 resamples of 0.5° scoring blocks), weighted by cell and by block as indicated in the key; the grey band marks ±0.5~cm.
(a) Smallest tested label count $n^{*}$ (line from the previous grid value; Methods) at which recalibration lowered error relative to the source-coefficient Stefan model (Recalibration) and residual ML relative to the recalibrated model (Residual ML); Not reached marks targets without $n^{*}$.
Rows follow the ten-label bias; sub-regions are not independent.
(b) Error change with all labels, relative to the source-coefficient Stefan model, of the method chosen by cross-validation within the target labels, against the absolute mean bias of that model at ten labels (Spearman $\rho = 0.38$; 95\% CI 0.17 to 0.55; 28 targets; signed bias, $\rho = -0.01$); the Tibetan Plateau lies outside the axes (bias 227.7~cm, error change −159.6~cm).
(c) Cross-validation selection minus a fixed residual recipe (residual weight 0.25), four main regions at 10 and all labels, the Lena Delta and Canada at 40 and 160 labels.
The selection was non-inferior (margin 0.5~cm) at 40 and 160 labels and had 1.3~cm lower error at 160 labels (95\% CI 0.7 to 1.7~cm); the 0.87~cm reduction at 40 labels and the all-label non-inferiority depend on the split-independence assumption, and non-inferiority did not hold at ten labels.
The change came from Canada (−2.1 and −2.7~cm; Lena Delta +0.35 and +0.07~cm).
(d) Error change from recalibration and from residual ML beyond recalibration with all labels, against the same bias; [XA: Fig 4d | 사후 분석(재현(비맹검)) | 2.1 결과 범주(양, 확인하지 못함, 음)와 CE1, CE2 판].
Panels b and c were designed after earlier results had been viewed; panel d is post hoc.

\medskip
% ---- Figure 5 -------------------------------------------------------------
% Sources: C6 README 1.3, 2.1, 2.2 (within-region strategy rows), 2.5 (transfer
% rows). Common-resampling verdicts read from the within-region and transfer
% placement rows of the source tables named in C6 README 2.0: Lena block
% stratification at 100 labels, Lena active selection at 200 labels and Canada
% covariate spread at 40 transfer labels depend on the split-independence
% assumption (not recorded in C6 README). The Alaskan sub-region rows of the spec
% draft are not plotted in panel d and are left to the Supplementary Information.
% Rendered Fig5 (fig5.py, Fig5_values.txt): maps a-c (Random, Block stratified, Covariate spread; no active-selection map,
% spec decision D-12 not approved), d = budget curves, e = transfer forest. Panel letters here follow the rendered figure.
\noindent Figure 5. Spreading labels across blocks and covariate space lowered error versus random cells in Canada, not in the Lena Delta.
(a to c) One draw of 100 labels in Canada (within-region test, split 1) by the three placement procedures named above the maps (Methods); small dots are candidate cells in the label blocks, and circle area is proportional to the labels selected per site (blocks closer than 40~km pooled).
The maps illustrate the procedures and do not rank sites.
(d) Error change relative to random cells for Anchor + residual ML (residual weight 0.25) against the number of labels within Canada, the Lena Delta and Alaska.
(e) The same contrast in transfer with ten and forty labels.
Intervals are 95\% block-bootstrap confidence intervals (CIs; 10,000 resamples of 0.5° scoring blocks, conditional on the label draws), weighted by cell and by block as indicated in the key; the grey band marks ±0.5~cm.
In panel d, both spreading strategies lowered error in Canada by 2.5--3.2~cm at 50 and 100 labels and by 1.3--1.4~cm at 200 labels.
In the Lena Delta, block stratification changed error by +1.5 to +2.0~cm at 20--500 labels, with higher error under both weightings at 100 labels (+1.7~cm; depends on the split-independence assumption).
In Alaska the weightings disagreed at 50--200 labels (block-equal −1.1 to −1.7~cm, cell-weighted +0.28 to +0.74~cm).
Variance-based active selection had 0.69 to 2.8~cm higher error at 50--200 labels in Alaska and the Lena Delta (one of six contrasts depends on the split-independence assumption) and smaller gains than spreading in Canada.
In panel e, covariate spread lowered error in Canada at forty labels by 2.4~cm (95\% CI 0.5 to 3.4~cm; depends on the split-independence assumption); at ten labels no region had higher error under both weightings, but the cell-weighted interval of block stratification in the Lena Delta lay above zero (+1.9~cm).
The placement tests were designed after earlier results had been viewed and registered before running.
Polar stereographic projection centred on Canada, scale bar in panel a; coastlines, Natural Earth 50~m; Cartopy [미확인: 판].

\medskip
% ---- Figure 6 -------------------------------------------------------------
% Sources: C1 README 1.3, 2.1 (learner rows; five learners with agreeing CIs),
% 2.5 (post hoc counts 18 and 2 of 30; 17 + 10 + 3 targets), 2.7 (year-matched
% baseline); abstract bundle row for the main CatBoost learner (2.2459, rounded
% 2.2). Interval resamples: c 10,000 (v2 figure meta), d 1000 with calibration
% fixed (interval plan; one-stage columns chosen in figure spec 8.3), e 10,000.
% Placeholder: spec section 11.
\noindent Figure 6. Without target labels, 2 of 30 targets lost over 2~cm with physics pseudo-labels and 18 with direct ML.
Intervals in panel c are 95\% block-bootstrap confidence intervals (CIs; 10,000 resamples of 0.5° scoring blocks), weighted by cell and by block as indicated in the key, with the grey band at ±0.5~cm; bars in panels d and e are cell-weighted 95\% CIs (d, 1000 resamples, calibration fixed; e, 10,000 resamples).
(a, b) Error change relative to the source-coefficient Stefan model for direct ML (a) and direct ML trained with physics pseudo-labels (b) at each held-out target whose source excludes the parent region (17 targets); shared colour scale.
The title counts are a post hoc description of 30 targets (adding 10 sub-regions held out within their region and 3 expanded-label editions; threshold on point estimates).
(c) Error change of direct ML by learner and of two physics baselines, four-region mean.
The random forest and four neural networks configured by leave-one-source-region-out cross-validation had 2.4--5.2~cm higher error.
The other five learners, including the main CatBoost learner (2.2~cm; 95\% CI 1.1 to 3.4~cm; Holm-adjusted $P = 0.023$), gave the same verdict under the primary intervals, but this verdict depends on the split-independence assumption, as does the 1.0~cm lower error of Year-matched Stefan.
[XG: Fig 6c | 결과 열람 뒤 설계, 등록 이탈(WRAPUP 10) | 2.7 XG-1 행]
(d) Coverage of 90\% prediction intervals against mean width without labels for five interval scalings, four-region mean and regional values; the line marks nominal 0.90.
(e) Coverage of anchor plus residual ML intervals with 40 and 160 labels against width relative to zero labels; the grey band marks 0.85 to 0.95.
Polar stereographic projection (45°~W; scale true at 70°~N) with a Lambert azimuthal equal-area Tibetan inset; coastlines, Natural Earth 50~m; permafrost zones [미확인: 자료 이름과 판]; Cartopy [미확인: 판].

\medskip
% ---- Figure 7 -------------------------------------------------------------
% Sources: C4 README 2.1 (Alaska rows; the 1000-label row agrees under common
% resampling, 500 and all do not), 2.2 (Lena 23-24 splits), 2.6 (floors, 19%);
% C8 README 2.1, 2.2 (about 99%: 1 - 0.0086), 2.4 (three-region within-cell row;
% common-resampling verdict undetermined in the source table); Canadian floor
% line = 17.90 - 28.55 and 17.90 - 28.66 cm (figure spec 9.6); panel b mask and
% 343 locations from figure spec 9.3. Placeholder: spec section 11.
\noindent Figure 7. Within Alaska, residual ML gains over recalibrated Stefan were below 1~cm and came from climate-grid-scale bias correction.
Intervals are 95\% block-bootstrap confidence intervals (CIs; 10,000 resamples of 0.5° scoring blocks), weighted by cell and by block as indicated in the key; the grey band marks ±0.5~cm.
(a) Within-region tests (half of the 0.5° blocks scored, up to 25 splits): error change of Anchor + residual ML (cross-validated residual weight) and Direct ML relative to the Stefan model recalibrated on the same labels; Error floor is the covariate-conditional floor minus the recalibrated Stefan RMSE, and the Canadian line (−10.8 to −10.6~cm) is off the axis.
In Alaska with 1000 labels, Anchor + residual ML had 0.54~cm lower error (95\% CI 0.32 to 0.69~cm), removing 19\% of the squared error above the floor; the reductions with 500 and all labels (0.39 and 0.52~cm) did not hold under common resampling, and no difference was established in the Lena Delta or Canada.
(b) Alaska, all labels: error change per 0.5° scoring block, with label locations; hatching marks blocks with fewer than 10 scoring cells summed over the 25 splits.
(c) Error change split between and within ERA5-Land climate cells; in Alaska about 99\% of the squared-error reduction lay between cells (within cells −0.01~cm, equivalent), and in the three-region mean the within-cell predictions increased error by 0.09~cm (depends on the split-independence assumption).
(d) Recommended method per label interval at its error change relative to the source-coefficient Stefan model (Lena Delta and Canada transfer mean; 320 and 1000 labels untested), with placement and diagnosis steps below; selection from 40 labels rests on non-inferiority to a fixed recipe (Fig.~4c).
(e) [XC: Fig 7e | 결과 열람 뒤 설계, 재사용 지역의 재검정(비맹검 부분 포함) | 2.3 XC-1, XC-2 포레스트].
Panels a to d were designed after viewing the transfer results.
Polar stereographic projection centred on Alaska (scale true at 70°~N); coastlines, Natural Earth 50~m; Cartopy [미확인: 판].

% ---- Table 1 caption ------------------------------------------------------
% Sources: figure spec 10.2 to 10.4; D README 2 (row counts, source coefficient, Alaska share
% 78-94%); cell index 0.009 degrees from the v2 Table 1 caption. If a cell of the
% assembled table has no value, write "n.a." and add "n.a., not available." here
% (style guide 3.1). The table body is generated separately (Table1_data.tex).
\section*{Table caption}

\noindent Table 1. Target regions, labels and data sources.
Labels is the number of label rows; 1~km locations groups the rows by a 0.009° cell index; 0.5° blocks are the units of the label and scoring split; Splits is the number of valid block splits in the transfer design; $E_0$ is the Stefan coefficient (cm~(°C~d)$^{-1/2}$) fitted on the source pool when the region is held out; Alaska share of source cells is the fraction of source cells located in Alaska.
Main regions were used in earlier experiments and are retested here; Alaska is a reference target that is not averaged with them; new regions were added after the label tables had been fixed (licence-verified edition).
Sub-regions, targets with point estimates only and the North Atlantic group are listed in Supplementary Table~S1.
Access conditions are given under Data availability.


\clearpage
% The class wraps 'table' in threeparttable, which cannot hold the boxed tabular; the original
% float (tableorg, saved by the class) is used so that the 8-column table can be centred across
% the text block (natural width about 14 mm wider than the 31 pc single-column text block).
\begin{tableorg}[p]
\caption{Target regions, labels and data sources.}\label{tab:data}\par\smallskip
\tablebodyfont\setlength{\tabcolsep}{4pt}%
\makebox[\textwidth][c]{\TableOneTabular}
\end{tableorg}
\clearpage

\PlaceFigure{1}{Fig1.pdf}
\PlaceFigure{2}{Fig2.pdf}
\PlaceFigure{3}{Fig3.pdf}
\PlaceFigure{4}{Fig4.pdf}
\PlaceFigure{5}{Fig5.pdf}
\PlaceFigure{6}{Fig6.pdf}
\PlaceFigure{7}{Fig7.pdf}

\end{document}
